\documentclass[11pt]{book}
\usepackage[utf8]{inputenc}
\usepackage[T1]{fontenc}
\usepackage{lmodern}
\usepackage{textcomp}
\usepackage{upquote}   % straight ' and ` in code, so examples can be pasted
\usepackage[margin=1.25in]{geometry}
\usepackage{verbatim}
\usepackage{makeidx}
\usepackage[colorlinks=true,linkcolor=black,urlcolor=blue,citecolor=black]{hyperref}
\makeindex

% Code examples a size smaller, so that 80 columns fit the line
\makeatletter
\renewcommand{\verbatim@font}{\normalfont\small\ttfamily}
\makeatother
\setlength{\emergencystretch}{2em}

% A straight quote for inline code, as in \code{\q a} or \code{\#\q}
\newcommand{\q}{\textquotesingle}
% Inline code
\newcommand{\code}[1]{\texttt{#1}}
% The tilde of a format directive, as in \code{\td a}
\newcommand{\td}{\textasciitilde}

\title{Understanding Common Lisp}
\author{\href{https://davecooper.name}{Dave Cooper}}
\date{2026 Edition}

\begin{document}

\frontmatter

\maketitle

\clearpage
\thispagestyle{empty}
\vspace*{\fill}
\begin{center}
Copyright \copyright{} 2003, 2026 Franz Inc. and
\href{https://davecooper.name}{Dave Cooper}.

\medskip
First published in 2003 as \emph{Basic Lisp Techniques}.

\medskip
This book is licensed under the Creative Commons
Attribution-ShareAlike 4.0 International License,\\
\url{https://creativecommons.org/licenses/by-sa/4.0/}.

\medskip
The source is at
\url{https://github.com/genworks/understanding-common-lisp}.
\end{center}
\vspace*{\fill}
\clearpage

% Preface
\chapter*{Preface to the 2026 Edition}
\addcontentsline{toc}{chapter}{Preface to the 2026 Edition}

Franz Inc. commissioned this book and published it in 2003 as
\emph{Basic Lisp Techniques}. It has been a free download from their
website ever since, and it has served many programmers as a quick
start in the language. The title it was meant to have was
\emph{Understanding Common Lisp}, and this edition takes that name.

The language has hardly changed since 2003. Common Lisp is defined by
a standard that was finished in 1994, and a program written to that
standard then still runs now. What has changed is everything around
the language. There are several free implementations that compile to
fast native code. A library manager puts a couple of thousand
open-source libraries one command away. The editor support that used
to come with one vendor's product works with every implementation.
And many of us now write Lisp alongside an AI agent that can evaluate
code in the same running image we are working in.

The first edition was written around one implementation, Allegro CL.
This one is written so that the examples run on any Common Lisp. They
are shown as SBCL prints them, since SBCL is free and is what most
newcomers install first, and every one of them was run before it went
into the book.

The book is still short. It is for a working programmer who wants to
read and write Common Lisp with some confidence in a few sittings, and
who will go on to the longer books and the reference material after
that. The first edition ended with a sample application. That chapter
is gone: there are plenty of real programs to read online now, and
Chapter~5 says where to find them.

Corrections and improvements are welcome at the address on the
copyright page.

My thanks to Franz Inc., who asked for the first edition and have kept
it available for more than twenty years.

\begin{flushright}
Dave Cooper\\
\url{https://davecooper.name}
\end{flushright}

% Foreword
\chapter*{Foreword to the 2003 Edition\footnote{This foreword was written by Franz Inc. for the first edition, and is reprinted as it appeared there.}}
\addcontentsline{toc}{chapter}{Foreword to the 2003 Edition}

Computers, and the software applications that power them, permeate every facet of our
daily lives. From groceries to airline reservations to dental appointments, our reliance
on technology is all-encompassing. And, it's not enough. Every day, our expectations of
technology and software increase:

\begin{itemize}
\item smart appliances that can be controlled via the internet
\item better search engines that generate information we actually want
\item voice-activated laptops
\item cars that know exactly where to go
\end{itemize}

The list is endless. Unfortunately, there is not an endless supply of programmers and
developers to satisfy our insatiable appetites for new features and gadgets. Every day,
hundreds of magazine and on-line articles focus on the time and people resources needed to
support future technological expectations. Further, the days of unlimited funding are over.
Investors want to see results, fast.

Common Lisp (CL) is one of the few languages and development options that can meet
these challenges. Powerful, flexible, changeable on the fly --- increasingly, CL is playing
a leading role in areas with complex problem-solving demands. Engineers in the fields of
bioinformatics, scheduling, data mining, document management, B2B, and E-commerce
have all turned to CL to complete their applications on time and within budget. CL,
however, no longer just appropriate for the most complex problems. Applications of modest
complexity, but with demanding needs for fast development cycles and customization, are
also ideal candidates for CL.

Other languages have tried to mimic CL, with limited success. Perl, Python, Java,
C++, C\# --- they all incorporate some of the features that give Lisp its power, but their
implementations tend to be brittle.

The purpose of this book is to showcase the features that make CL so much better
than these imitators, and to give you a ``quick-start'' guide for using Common Lisp as a
development environment. If you are an experienced programmer in languages other than
Lisp, this guide gives you all the tools you need to begin writing Lisp applications. If you've
had some exposure to Lisp in the past, this guide will help refresh those memories and shed
some new light on CL for you.

But be careful, Lisp can be addicting! This is why many Fortune 500 companies will
use nothing else on their 24/7, cutting-edge, mission-critical applications. After reading
this book, trying our software, and experiencing a 3 to 10 times increase in productivity,
we believe you will feel the same way.

\tableofcontents

\mainmatter

% Chapter 1
\chapter{Introduction}

\section{The Past, Present, and Future of Common Lisp\index{Common Lisp}}

\subsection{Lisp Yesterday}

John McCarthy\index{McCarthy, John} worked out the basic principles of Lisp in 1958, at
MIT. The name stood for ``LISt Processing'': the language was built
for manipulating symbols and lists of symbols, where the other
languages of the day were built for arithmetic. A list, in this
sense, is just a sequence of items, like a shopping list or a
checklist.

The idea that set Lisp apart was that a Lisp program is itself written
as lists. The language can therefore read, build and transform
programs with the same operations it uses on any other data. Everything
that makes Lisp unusual follows from that.

Lisp spread quickly through the research world and split into many
dialects. In the early 1980s the people behind the main ones agreed to
merge them into one language, Common Lisp. Guy Steele described it in
the book \emph{Common Lisp: the Language} in 1984, and after ten more
years of work it became an ANSI standard in 1994. It was the first
object-oriented language to be standardized by ANSI\index{ANSI standard}.

\subsection{Lisp Today}

The standard has not changed since 1994, and this turns out to be one
of the best things about the language. A Common Lisp program written
thirty years ago still runs today, usually without touching a line of
it, on implementations and machines that did not exist when it was
written. Libraries do not rot when a new version of the language
comes out, because one does not come out. Code you write this year
will not need rewriting for the language's sake.

The language can stand still because it can be extended from the
inside. Most of what other languages add in new versions (a new kind
of loop, pattern matching, a new way to manage resources) can be added
to Common Lisp by a library, and looks no different from what is built
in. Section~\ref{sec:macros} shows how.

Meanwhile the implementations and tools have kept moving. There are
several maintained implementations of Common Lisp, most of them free
and open source, and the best of them compile to efficient machine
code. Some two thousand
open-source libraries can be installed with one command. Common Lisp
is in daily commercial use in engineering design, scheduling and
logistics, finance, and research computing, and it has a small, steady
community of people who chose it on purpose.

\subsection{Lisp Tomorrow}

For most of the history of software, people have hoped for tools that
would let them describe a program at a high level and have the detailed
code written for them. Lisp's answer has always been that the
high-level description and the program should be the same thing,
written in one language. A Common Lisp program can generate other
Common Lisp code, normally at compile time, so you can build the
language up toward your problem until the program reads like a
description of it.

That matters more now, not less. A growing share of code is drafted by
AI models and then checked by people. A language with a small, regular
syntax, a precise standard, and a running image that can be asked what
is true right now is a good place to do that kind of work: the model
can try an expression and see the answer before it proposes anything,
and so can you. Section~\ref{sec:agents} says more.

Several long-standing features of Common Lisp also save time and
people:

\begin{itemize}
\item \textbf{Automatic memory management\index{garbage collection}.} You do not
allocate and free memory by hand. Lisp was the first language to have
a garbage collector.

\item \textbf{Dynamic typing\index{dynamic typing}.} Values have types, but variables
do not have to. You can write a function without declaring the type
of anything, and add declarations later where they help the compiler
or the reader.

\item \textbf{Redefinition while running.} You can add and replace
functions, classes and data in a running program without stopping it.
This is how Lisp is normally developed, and it also means that a
server can be inspected and patched while it stays up.

\item \textbf{A large standard language.} Hash tables, arbitrary-precision
integers, exact fractions, multi-dimensional arrays, a full object
system, formatted output, and a condition system for handling errors
are all part of the standard.

\item \textbf{Portability.} Because the standard is precise, code moves
between implementations and operating systems with little or no
change.
\end{itemize}

\section{Is Lisp Slow? Is It Big?}

The old complaints about Lisp were that it was too slow and needed too
much memory. They were fair on the machines of the 1970s. They have
not been true for a long time. Compiled Common Lisp runs at the speed
of other compiled languages, and a Lisp image of forty megabytes is
small next to a web browser tab.

The cost that matters today is the time of the people writing the
program. Common Lisp teams tend to be small because one person with a
good grasp of the language can build and maintain a great deal.

\section{The CL Model of Computation}

In most languages you work in a cycle: edit the source, build, run the
program from the beginning, and see what happens. In Common Lisp you
start the Lisp once and leave it running. You type an expression and
it is evaluated at once. You change a function and send just that
function to the running Lisp, which replaces the old definition and
carries on, with all its data still in memory. You test the change in
the next second.

This running Lisp is called the \emph{image}\index{image}, and the prompt where
you talk to it is called the \emph{REPL}\index{REPL}, for read-eval-print loop: it
reads an expression, evaluates it, prints the result, and waits for
the next one. Working this way changes how you write programs. You
build them in small pieces and try each piece as you go, so you find
out early when you have misunderstood something.

The next chapter shows how to set this up.


% Chapter 2
\chapter{Operating a CL Development Environment}

This chapter covers the hands-on side of working with Common Lisp:
installing it, talking to it, connecting it to your editor, loading
libraries, and finding out what went wrong when something does. If
Lisp code is completely new to you, you may want to glance ahead at
Chapter~3 as you read.

It helps to think of a running Common Lisp as a small operating system
of its own, sitting on top of whichever operating system your machine
runs. You start it, and then you load things into it, ask it
questions, and change it while it runs. People once built whole
workstations this way, with Lisp all the way down.

Have a Lisp running while you read, and type the examples in.

\section{Choosing an Implementation}

Common Lisp is a standard with several implementations\index{implementations}, in the way
that C and SQL are. The ones you are most likely to meet:

\begin{itemize}
\item \textbf{SBCL}\index{SBCL} (Steel Bank Common Lisp), \url{https://www.sbcl.org}.
Free and open source, with a compiler that produces fast native code
and gives thorough warnings. The most widely used implementation.

\item \textbf{CCL}\index{CCL} (Clozure CL), \url{https://ccl.clozure.com}. Free and
open source. Compiles very quickly, and is pleasant to develop in.

\item \textbf{ECL}\index{ECL} (Embeddable Common Lisp),
\url{https://ecl.common-lisp.dev}. Free. Compiles by way of C, so
it can be linked into C and C++ programs and runs almost anywhere.

\item \textbf{ABCL}\index{ABCL} (Armed Bear Common Lisp), \url{https://abcl.org}.
Free. Runs on the Java virtual machine and can call Java libraries.

\item \textbf{Allegro CL}\index{Allegro CL}, from Franz Inc., \url{https://franz.com}, and
\textbf{LispWorks}\index{LispWorks}, \url{https://www.lispworks.com}. Commercial
products with their own development environments, graphical toolkits
and support, each with a free edition for learning.
\end{itemize}

They all implement the same language, and what you learn on one
carries over to the others. This book uses SBCL. Where something in
the text is specific to SBCL, it says so.

\section{Installing a CL Environment}

SBCL is packaged for the common systems:

\begin{verbatim}
sudo apt install sbcl         (Debian, Ubuntu)
brew install sbcl             (macOS with Homebrew)
\end{verbatim}

\noindent
For Windows there is an installer at \url{https://www.sbcl.org}. The
packaged versions can be a few releases behind; the website has
current binaries if you want them.

If you use Docker\index{Docker}, you can try an implementation without installing
anything. The \code{clfoundation} images on Docker Hub cover most of
the free implementations:

\begin{verbatim}
docker run -it --rm clfoundation/sbcl
\end{verbatim}

\section{Running CL in a Terminal}

The simplest way to run CL is from a terminal. Few people work this
way all day, because an editor connected to Lisp is far more
convenient, but it is the place to start, and it is still useful for
logging in to a Lisp that is running on a remote machine.

\subsection{Starting and Stopping}

To start SBCL, type

\begin{verbatim}
sbcl
\end{verbatim}

\noindent
You will see a few lines of introduction and then a prompt:

\begin{verbatim}
This is SBCL 2.6.9, an implementation of ANSI Common Lisp.
More information about SBCL is available at <http://www.sbcl.org/>.
...
*
\end{verbatim}

CL has entered its read-eval-print loop\index{REPL} and is waiting for you to type
something. It will read what you type, evaluate it to get a return
value, and print that value:

\begin{verbatim}
* (+ 2 2)
4
\end{verbatim}

To leave SBCL, type

\begin{verbatim}
(exit)
\end{verbatim}

\noindent
Each implementation has its own name for this (in CCL it is
\code{(quit)}). Typing the end-of-file character, Control-D, works in
most of them.

The bare terminal prompt has no command history and no line editing.
If you plan to spend any time there, install \code{rlwrap} and start
Lisp as \code{rlwrap sbcl}.

\subsection{Compiling and Loading a File}

\textbf{Try it:} With your text editor, create a file
\code{/tmp/hello.lisp} containing the following (do not worry about
understanding it yet):

\begin{verbatim}
(defun hello ()
  (write-string "Hello, World!"))
\end{verbatim}

Now, at the CL prompt, compile and load the file. The text after each
\code{*} prompt is what you type; everything else is printed by CL:

\begin{verbatim}
* (compile-file "/tmp/hello.lisp")
; compiling file "/tmp/hello.lisp" (written 05 OCT 2026 06:48:14 PM):

; wrote /tmp/hello.fasl
; compilation finished in 0:00:00.000
#P"/tmp/hello.fasl"
NIL
NIL
* (load "/tmp/hello")
T
\end{verbatim}

\code{Compile-file}\index{compile-file@\texttt{compile-file}} translates the source file into a binary file of
machine code. The binary is called a \emph{fasl} file\index{fasl file}, for ``fast
loading.'' \code{Load}\index{load@\texttt{load}} brings a file's definitions into the running
Lisp. Given a name with no extension, it picks the fasl file if there
is one.

You can also \code{load} a source file directly, and the effect is the
same, except that it has to be processed each time. Fasl files are
specific to the implementation, its version, and the kind of machine.
You do not pass them around; each Lisp makes its own from the source.

Now call the function you defined:

\begin{verbatim}
* (hello)
Hello, World!
"Hello, World!"
\end{verbatim}

The text appears twice. The first line is the function printing the
string, as you asked it to. The second is the REPL printing the value
that the function returned, which happens to be the same string; it
is shown in double quotes because the REPL prints values in the form
you would type them. Chapter~3 has more on this.

The point for now is the way of working: definitions live in files,
you compile and load them into a running Lisp, and you call them from
the prompt.

\section{Running CL inside an Editor}

The way most people develop in CL is with the Lisp running as a
subprocess of their text editor. The editor then knows about the
running Lisp. With a keystroke you can send it the function you are
editing, ask it for a function's arguments or documentation, jump to
the source of anything that is defined, and look inside live data.

\subsection{Emacs with SLIME or Sly}

The traditional editor for Lisp is GNU Emacs\index{Emacs}, which is itself largely
written in a dialect of Lisp, and the standard connection between
Emacs and Common Lisp is SLIME\index{SLIME} (\url{https://slime.common-lisp.dev}).
Sly\index{Sly} is a fork of SLIME with some added conveniences. Both work with
every implementation listed above, and what follows applies to both.

If you are new to Emacs, run its built-in tutorial first (it is under
the Help menu). Emacs commands are written as key chords: \code{C-x}
means hold Control and press x; \code{M-x} means hold Meta, which is
the Alt or Option key, and press x.

To install SLIME, add the MELPA package archive to Emacs if you have
not already (\url{https://melpa.org} has the two lines to copy), then
type \code{M-x package-install}, press Return, and give the name
\code{slime}. Then put this line in your Emacs init file to say
which Lisp to run:

\begin{verbatim}
(setq inferior-lisp-program "sbcl")
\end{verbatim}

Now \code{M-x slime} starts Lisp and opens a REPL in an Emacs window.
The prompt shows the current package, a notion explained in
Section~\ref{sec:packages}:

\begin{verbatim}
CL-USER> (+ 2 2)
4
\end{verbatim}

\noindent
The rest of this book shows examples with this prompt.

Table~\ref{tab:slime} lists the commands you will use most.

\begin{table}[ht]
\centering
\begin{tabular}{ll}
\hline
\textbf{Key} & \textbf{Action} \\
\hline
\code{C-c C-c} & Compile the definition the cursor is in \\
\code{C-c C-k} & Compile and load the whole file \\
\code{C-x C-e} & Evaluate the expression before the cursor \\
\code{C-c C-z} & Switch to the REPL \\
\code{M-.}     & Go to the definition of the name at the cursor \\
\code{M-,}     & Come back from there \\
\code{C-c C-d d} & Describe the name at the cursor \\
\code{C-c RET} & Expand the macro call at the cursor \\
\code{M-p}, \code{M-n} & In the REPL: previous and next input \\
\code{C-c C-c} & In the REPL: interrupt a running computation \\
\hline
\end{tabular}
\caption{Common SLIME and Sly commands}
\label{tab:slime}
\end{table}

While you type a call to a function, its argument list is shown at the
bottom of the window. Emacs indents Lisp code for you and shows which
parentheses match, so you read the structure of the code from its
indentation and never count parentheses.

\textbf{Try it:} Open \code{/tmp/hello.lisp} in Emacs, with SLIME
running. Put the cursor inside the definition of \code{hello} and
press \code{C-c C-c}. Switch to the REPL with \code{C-c C-z} and type
\code{(hello)}. Now go back to the file, change the string to
\code{"Hello, CL World!"}, press \code{C-c C-c} again, and call
\code{(hello)} again in the REPL. You have changed a running program
without restarting it.

\subsection{Other Editors}

Emacs is not required. Visual Studio Code has the Alive\index{Alive} extension,
which gives a REPL, evaluation from the editor, a debugger and an
inspector. Vim and Neovim have Vlime and Slimv. Lem\index{Lem}
(\url{https://lem-project.github.io}) is an Emacs-like editor written
in Common Lisp, which works with Lisp as soon as it starts. Allegro
CL and LispWorks each come with their own environment.

Whatever you choose, make sure it gives you the essentials: a REPL,
a key to send the current definition to the running Lisp, and
automatic indentation.

\section{The Init File}

When CL starts, it looks for an initialization file in your home
directory and loads it\index{init file}. In SBCL this is \code{.sbclrc}; in CCL it is
\code{.ccl-init.lisp}. The file holds ordinary Lisp, and is the place
to load the tools you want in every session. The next section puts
something in it.

\section{Libraries: Quicklisp and ASDF}

\subsection{Loading Libraries with Quicklisp}

Quicklisp\index{Quicklisp} (\url{https://www.quicklisp.org}) downloads and loads
open-source libraries, along with whatever they depend on. To install
it, download one file and load it:

\begin{verbatim}
curl -O https://beta.quicklisp.org/quicklisp.lisp
sbcl --load quicklisp.lisp
\end{verbatim}

\noindent
and then, at the Lisp prompt:

\begin{verbatim}
* (quicklisp-quickstart:install)
* (ql:add-to-init-file)
\end{verbatim}

The second form adds a few lines to your init file so that Quicklisp
is there whenever you start Lisp. From then on, one call fetches a
library and loads it:

\begin{verbatim}
CL-USER> (ql:quickload "alexandria")
To load "alexandria":
  Load 1 ASDF system:
    alexandria
; Loading "alexandria"

("alexandria")
CL-USER> (alexandria:flatten '(1 (2 (3 4)) 5))
(1 2 3 4 5)
\end{verbatim}

Alexandria, loaded here, is a collection of small utilities that most
other libraries use. \code{(ql:system-apropos "json")} searches the
library list by name. Chapter~4 uses several libraries, and
Chapter~5 says where to look for more.

Quicklisp is the common starting point, and there are other tools for
the same job, such as ocicl and Qlot, which add things like pinning
the exact versions a project uses.

\subsection{Defining Your Own System with ASDF}
\label{sec:asdf}

A program of any size is spread over several files, some of which have
to be loaded before others. ASDF\index{ASDF} is the standard tool for describing
that. It comes with every implementation mentioned in this chapter,
and it is what Quicklisp uses underneath. In ASDF, a library or
program is called a \emph{system}\index{system}, and is described by a file
ending in \code{.asd}.

Here is a complete small system, in a directory named
\code{greeter}. First \code{greeter.asd}:

\begin{verbatim}
(defsystem "greeter"
  :description "A very small example system."
  :version "0.1.0"
  :depends-on ("alexandria")
  :components ((:file "package")
               (:file "greet" :depends-on ("package"))))
\end{verbatim}

\noindent
then \code{package.lisp}:

\begin{verbatim}
(defpackage :greeter
  (:use :common-lisp)
  (:export #:greet))
\end{verbatim}

\noindent
and \code{greet.lisp}:

\begin{verbatim}
(in-package :greeter)

(defun greet (name)
  (format nil "Hello, ~a!" (string-capitalize name)))
\end{verbatim}

Put the \code{greeter} directory inside \code{\textasciitilde/common-lisp/},
which is where ASDF looks by default, and you can load the system by
name like any other library:

\begin{verbatim}
CL-USER> (ql:quickload "greeter")
To load "greeter":
  Load 1 ASDF system:
    greeter
; Loading "greeter"

("greeter")
CL-USER> (greeter:greet "reader")
"Hello, Reader!"
\end{verbatim}

ASDF compiles each file the first time and whenever it has changed,
keeps the fasl files out of your source directory, and loads the files
in an order that respects the dependencies you declared. Quicklisp
fetches Alexandria if it is not yet installed.

Start every project this way, however small. It costs three short
files, and it means anyone, including you a year from now, can load
the program with one line.

\section{Scripts and Standalone Programs}

\subsection{Running a File as a Script}

CL is normally used interactively, but it will also run a file from
the command line and exit\index{scripting}. Here is \code{script.lisp}, which copies
its input to its output with leading spaces removed from each line:

\begin{verbatim}
;;; Print each line of standard input with leading spaces removed.
(loop for line = (read-line *standard-input* nil)
      while line
      do (write-line (string-left-trim " " line)))
\end{verbatim}

\begin{verbatim}
sbcl --script script.lisp < datafile1 > datafile2
\end{verbatim}

A semicolon starts a comment\index{comments}, which runs to the end of the line. By
convention one semicolon is used for a comment after code on the same
line, two for a comment on its own line inside a definition, and three
for a comment at the top level of a file.

With \code{--script}, SBCL skips your init file, so Quicklisp is not
loaded and neither is ASDF. A script that needs them starts by
loading them itself.

\subsection{Building an Executable}

Because a running Lisp is an image\index{image} in memory, most implementations
can write that image to disk, and can write it as a program that
starts in a function of your choosing. This is how a Lisp application
is delivered to someone who does not have Lisp installed. In SBCL:

\begin{verbatim}
(defun main ()
  (format t "Hello from a standalone program!~%"))

(sb-ext:save-lisp-and-die "hello-app"
                          :toplevel #'main
                          :executable t)
\end{verbatim}

This writes a file \code{hello-app} that runs on any machine of the
same kind. It is about forty megabytes, because it contains the whole
Lisp, compiler included. Each implementation has its own function
for this. If you need one build recipe for several implementations,
ASDF can do it; look up \code{program-op} in its manual.

\section{Debugging in CL}
\label{sec:debugging}

Common Lisp has one of the best debugging and error-handling
arrangements of any language, and it is always on. There is no
separate debug mode and no debug build.

\subsection{The Debugger}

When an error happens in most languages, the program unwinds to the
top and dies, and you read a stack trace to work out what the state
must have been. When an error happens in Common Lisp, the program
\emph{stops where the error happened}, with everything still in place,
and you are put into the debugger\index{debugger}:

\begin{verbatim}
* (defun oops (x) (/ 10 x))
OOPS
* (oops 0)

debugger invoked on a DIVISION-BY-ZERO in thread
#<THREAD tid=59 "main thread" RUNNING {1200038003}>:
  arithmetic error DIVISION-BY-ZERO signalled
Operation was (/ 10 0).

Type HELP for debugger help, or (SB-EXT:EXIT) to exit from SBCL.

restarts (invokable by number or by possibly-abbreviated name):
  0: [ABORT] Exit debugger, returning to top level.

(SB-KERNEL::INTEGER-/-INTEGER 10 0)
0]
\end{verbatim}

The prompt \code{0]} is the debugger's. It is still a Lisp prompt:
you can evaluate any expression there, look at variables, and even
redefine the function that failed. \code{backtrace} shows the chain
of calls that led here. Typing \code{0}, or \code{abort}, takes the
restart numbered 0 and returns you to the normal prompt.

In SLIME or Sly the same information appears in its own window. Press
the number of a restart to choose it, \code{q} to abort back to the
REPL, Return on a line of the backtrace to see that call's local
variables, and \code{v} to jump to its source.

\subsection{Restarts}

The list headed ``restarts'' is the unusual part\index{restarts}. A restart is a
way of carrying on that the program itself has offered. Try opening
a file that does not exist:

\begin{verbatim}
* (with-open-file (in "/tmp/nope.conf") (read in))

debugger invoked on a FILE-DOES-NOT-EXIST in thread
#<THREAD tid=68 "main thread" RUNNING {1200038003}>:
  The file #P"/tmp/nope.conf" does not exist: No such file or directory

restarts (invokable by number or by possibly-abbreviated name):
  0: [CREATE   ] Reopen with :if-does-not-exist :create
  1: [CONTINUE ] Retry opening.
  2: [USE-VALUE] Try opening a different file.
  3: [ABORT    ] Exit debugger, returning to top level.
\end{verbatim}

You could now create the file from another window and choose restart
1, and the program would carry on from the middle as if nothing had
gone wrong. When the failure comes two hours into a long computation,
that is worth a great deal. Section~\ref{sec:conditions} shows how a
program handles errors itself, without the debugger.

\subsection{Stopping a Program Yourself}

To interrupt a computation that is taking too long, press Control-C in
the terminal, or \code{C-c C-c} in the SLIME REPL. You land in the
debugger, where you can look around and then either continue or abort.

To make a program stop at a particular place, call \code{break}\index{break@\texttt{break}} there:

\begin{verbatim}
(defun risky (n)
  (break "n is ~a" n)
  (* n 2))
\end{verbatim}

\noindent
Calling \code{(risky 21)} enters the debugger with the message ``n is
21'' and a \code{CONTINUE} restart, which resumes the function and
returns 42.

\subsection{Trace, Describe and Friends}

\code{Trace}\index{trace@\texttt{trace}} prints every call to a function and what it returned,
without your editing the function. It is especially good for seeing
how a recursive function works:

\begin{verbatim}
CL-USER> (defun our-length (list)
           (if (null list)
               0
               (+ (our-length (rest list)) 1)))
OUR-LENGTH
CL-USER> (trace our-length)
(OUR-LENGTH)
CL-USER> (our-length '(a b))
  0: (OUR-LENGTH (A B))
    1: (OUR-LENGTH (B))
      2: (OUR-LENGTH NIL)
      2: OUR-LENGTH returned 0
    1: OUR-LENGTH returned 1
  0: OUR-LENGTH returned 2
2
CL-USER> (untrace our-length)
T
\end{verbatim}

The running Lisp can also tell you about itself.
\code{Describe}\index{describe@\texttt{describe}} prints what is known about any object,
\code{documentation}\index{documentation@\texttt{documentation}} returns a documentation string, and
\code{apropos}\index{apropos@\texttt{apropos}} finds names containing a piece of text:

\begin{verbatim}
CL-USER> (documentation 'mapcar 'function)
"Apply FUNCTION to successive tuples of elements of LIST and MORE-LISTS.
Return list of FUNCTION return values."
CL-USER> (apropos "HASH-TABLE-" :cl)
HASH-TABLE-COUNT (fbound)
HASH-TABLE-P (fbound)
HASH-TABLE-REHASH-SIZE (fbound)
HASH-TABLE-REHASH-THRESHOLD (fbound)
HASH-TABLE-SIZE (fbound)
HASH-TABLE-TEST (fbound)
WITH-HASH-TABLE-ITERATOR (fbound)
\end{verbatim}

\subsection{Timing}

Wrap any expression in the \code{time}\index{time@\texttt{time}} macro to see how long it takes
and how much memory it allocates. The expression's value is returned
as usual:

\begin{verbatim}
CL-USER> (defun fib (n)
           (if (< n 2)
               n
               (+ (fib (- n 1)) (fib (- n 2)))))
FIB
CL-USER> (time (fib 30))
Evaluation took:
  0.023 seconds of real time
  0.021844 seconds of total run time (0.021844 user, 0.000000 system)
  95.65% CPU
  65,398,250 processor cycles
  0 bytes consed

832040
\end{verbatim}

``Consed'' means allocated; Section~\ref{sec:cons} explains the word.
For finding out where a larger program spends its time, the
implementations provide profilers. SBCL's is called
\code{sb-sprof}.

\section{Working with an AI Agent}
\label{sec:agents}

Many programmers now work with an AI coding agent\index{AI agents}: a language model
that reads the project, proposes changes, and runs commands. Lisp
suits this way of working unusually well, for the same reasons it
suits a person at a REPL.

An agent working in most languages writes code, runs the whole program
or its tests, and reads the output. An agent connected to a running
Lisp image can do what you do: evaluate one expression and look at the
answer, ask the image for a function's argument list or documentation
instead of guessing, redefine one function and try it, and inspect the
state left behind by an error. Mistakes are caught an expression at a
time. The connection is usually made with the Model Context
Protocol\index{Model Context Protocol} (MCP), a common way of giving a model tools to call. One
open-source bridge between MCP and a running Lisp is lisply-mcp,
\url{https://github.com/genworks/lisply-mcp}.

Some advice that has held up:

\begin{itemize}
\item Have the agent evaluate what it writes, in the image, before it
shows it to you. Code that has only been written is a guess.

\item Read what it produces. To check an agent's Lisp you need what
this book teaches: what gets evaluated and when, what a function
returns, what is modified in place.

\item Let it work in a Lisp that is safe to break, such as a
container, and keep your own work in version control.
\end{itemize}

An agent makes it more useful to understand the language, not less.
You are the one who decides whether the program is right.


% Chapter 3
\chapter{The CL Language}

This chapter is a condensed tour of the language. It covers what you
need in order to read and write ordinary Common Lisp, and it leaves a
good deal out. Chapter~5 lists longer books and the reference
material for when you want the rest.

Type the examples in as you go. The text after the \code{CL-USER>}
prompt is what you type, and the lines after it are what Lisp prints.

\section{Overview of CL and its Syntax}

The first thing to notice about CL is that it uses prefix
notation\index{prefix notation}. The most common thing to do, in a program or at the
REPL, is to call a function, and you do that by writing a list: the
name of the function first, then its arguments, all inside
parentheses.

\begin{verbatim}
CL-USER> (+ 2 2)
4
\end{verbatim}

This expression is the function named by the symbol \code{+},
followed by the arguments 2 and 2. Evaluating it returns 4.

\subsection{How Expressions Are Evaluated}

When CL is asked to evaluate\index{evaluation} an expression, it follows a few rules:

\begin{enumerate}
\item A number evaluates to itself:

\begin{verbatim}
CL-USER> 99
99
\end{verbatim}

\item A string, which is text in double quotes, evaluates to itself:

\begin{verbatim}
CL-USER> "Be braver -- you can't cross a chasm in two small jumps."
"Be braver -- you can't cross a chasm in two small jumps."
\end{verbatim}

\item A symbol, which is a bare name, is treated as a variable and
evaluates to the variable's value. Variables come in
Section~\ref{sec:variables}.

\item A list is treated as a call. The first element names a function
(or a macro, which we will come to). The remaining elements are the
arguments. Each argument is evaluated, left to right, and the
function is then called with the resulting values:

\begin{verbatim}
CL-USER> (expt 2 5)
32
\end{verbatim}
\end{enumerate}

Since the arguments are themselves expressions, calls nest to any
depth:

\begin{verbatim}
CL-USER> (+ (+ 2 2) (+ 3 3))
10
\end{verbatim}

\textbf{Try it:} Type these at your prompt and convince yourself that
each result makes sense:

\begin{verbatim}
(+ 2)

(+)

(user-homedir-pathname)

(get-universal-time)    ; seconds since the start of the year 1900

(search "Dr." "I am Dr. Strangelove")

(subseq "I am Dr. Strangelove"
        (search "Dr." "I am Dr. Strangelove"))
\end{verbatim}

Do not worry about getting lost among the parentheses. Your editor
indents the code to show its structure and shows you which parentheses
match, and after a short while you read Lisp by its indentation, as
you read any other language.

\subsection{Lisp Syntax Simplicity}

What you have just read is nearly the whole syntax of the language.
The one thing to add is that some operators are \emph{macros}\index{macros} or
\emph{special operators}\index{special operators} instead of functions. They are written the
same way, with the operator's name first, but they do not necessarily
evaluate all their arguments; each has its own rule. You will meet
the common ones in this chapter and their rules are easy to remember.

If you come from a language such as C, Java, JavaScript or Python, you
have been using a mixture of notations: infix for arithmetic
(\code{2 + 2}), postfix for a few things (\code{i++}), and prefix for
everything else (\code{max(a, b)}). Lisp uses the prefix form for
everything and moves the opening parenthesis to the left of the
function's name. That is the whole difference. In return there are
no precedence rules to remember, and \code{+} can take any number of
arguments.

\subsection{Turning Off Evaluation}
\label{sec:quote}

Sometimes you want an expression itself and not its value. The
special operator \code{quote}\index{quote@\texttt{quote}} returns its argument without evaluating
it:

\begin{verbatim}
CL-USER> (quote a)
A
CL-USER> (quote (+ 1 2))
(+ 1 2)
\end{verbatim}

The first returns the symbol \code{A} itself, where a bare \code{a}
would have been looked up as a variable. The second returns a list of
three items, where the unquoted form would have been a call returning
3.

This is needed so often that there is an abbreviation: a single quote
character in front of an expression.

\begin{verbatim}
CL-USER> 'a
A
CL-USER> '(+ 1 2)
(+ 1 2)
\end{verbatim}

Notice that Lisp printed the symbol in upper case. By default the
reader converts the names of symbols to upper case as it reads them,
so \code{a}, \code{A}, \code{hello} and \code{Hello} in your source
all name symbols with upper-case names. People write their code in
lower case and Lisp answers in capitals.

\subsection{Fundamental CL Data Types}

Common Lisp has the data types you know from other languages, such as
numbers, strings and arrays, and some that may be new, such as
symbols and lists. This section introduces numbers, strings, symbols
and lists. The others come later in the chapter.

CL is dynamically typed\index{dynamic typing}: values have types, and variables can
hold a value of any type. The type of a value is checked when the
program runs. You may declare types where you want the compiler to
check or optimize something, and you never have to.

\subsubsection{Numbers}

Numbers\index{numbers} in CL form a family of types: integers, ratios, floating-point
numbers and complex numbers. Most of the time you just think
``number.'' The arithmetic functions take any mixture and return a
result of the appropriate kind.

\begin{verbatim}
CL-USER> (/ 1 3)
1/3
CL-USER> (+ 1/3 2/3)
1
CL-USER> (/ 1.0 3)
0.33333334
CL-USER> (expt 2 100)
1267650600228229401496703205376
CL-USER> (sqrt -4)
#C(0.0 2.0)
\end{verbatim}

Dividing two integers gives an exact ratio\index{ratios}, so one third is really
one third, and does not get rounded until you ask for a floating-point
number. Integers have no maximum size, so arithmetic on them never
overflows.

\subsubsection{Strings}

A string\index{strings} is a sequence of characters written between double quotes.
Strings can hold any Unicode text in all current implementations.

\begin{verbatim}
CL-USER> (string-upcase "hello")
"HELLO"
CL-USER> (concatenate 'string "Hello, " "World")
"Hello, World"
\end{verbatim}

A single character is written with the prefix \code{\#\textbackslash}, as in
\code{\#\textbackslash a} for the letter a and \code{\#\textbackslash Space} for a space.

\subsubsection{Symbols}

Symbols\index{symbols} are names. You have already used some: \code{+},
\code{expt} and \code{search} are symbols that name functions.
Symbols are also data in their own right. A program can hold a
symbol in a variable, put it in a list, and compare it with another
symbol, which is very fast because two occurrences of the same name
are the very same object.

A symbol can name a function and a variable at the same time without
any conflict. \code{List}, for example, is a built-in function, and
you are free to use \code{list} as a variable name as well. There is
no need to shorten it to \code{lst}.

To get a symbol itself where Lisp would otherwise evaluate it, quote
it, as in Section~\ref{sec:quote}:

\begin{verbatim}
CL-USER> 'my-symbol
MY-SYMBOL
\end{verbatim}

Symbols whose names begin with a colon are \emph{keywords}\index{keywords}. A keyword
evaluates to itself, so it needs no quote. Keywords are used as
labels and option names throughout the language:

\begin{verbatim}
CL-USER> :red
:RED
\end{verbatim}

\subsubsection{Lists}

A list\index{lists} is written as items between parentheses. The items can be
of any type, including other lists. Since an unquoted list is a call,
you quote a list when you mean it as data:

\begin{verbatim}
CL-USER> '(a b c)
(A B C)
CL-USER> '(a (b c) d)
(A (B C) D)
\end{verbatim}

The quote covers everything inside the list.

Now look at this list:

\begin{verbatim}
(defun hello () (write-string "Hello, World!"))
\end{verbatim}

It is also a program: the function definition from Chapter~2. Take a
moment to see that it meets the description of a list. It has four
items. The first two are symbols, the third is an empty list, and the
fourth is a list of a symbol and a string.

Every CL program is made of lists like this one. That is why it is
straightforward for a CL program to read, build and transform other CL
programs, and it is the main thing that sets Lisp apart from other
languages.

Besides quoting, you can make a list by calling the function
\code{list}\index{list@\texttt{list}}, which evaluates its arguments, as every function does, and
returns a list of the results:

\begin{verbatim}
CL-USER> (list 'a 'b (+ 2 2))
(A B 4)
\end{verbatim}

\subsection{Functions}

You define a named function with \code{defun}\index{defun@\texttt{defun}}. It takes a name, a list
of parameters, and a body of one or more expressions. The value of
the last expression in the body is the value the function returns;
there is no \code{return} statement.

\begin{verbatim}
CL-USER> (defun my-first-lisp-function ()
           (list 'hello 'world))
MY-FIRST-LISP-FUNCTION
CL-USER> (my-first-lisp-function)
(HELLO WORLD)
\end{verbatim}

\code{Defun} is a macro, so it does not evaluate its arguments in the
ordinary way: the name and the parameter list are taken as written,
with no quoting.

If the first thing in the body is a string, it is kept as the
function's documentation\index{documentation strings}:

\begin{verbatim}
CL-USER> (defun square (x)
           "Return X multiplied by itself."
           (* x x))
SQUARE
CL-USER> (square 4)
16
CL-USER> (documentation 'square 'function)
"Return X multiplied by itself."
\end{verbatim}

Names in Lisp are conventionally written in lower case with hyphens
between the words.

\subsection{Global and Local Variables}
\label{sec:variables}

\subsubsection{Global Variables}

Global variables\index{global variables} are established with \code{defparameter}\index{defparameter@\texttt{defparameter}} or
\code{defvar}\index{defvar@\texttt{defvar}}:

\begin{verbatim}
CL-USER> (defparameter *todays-temp* 90)
*TODAYS-TEMP*
CL-USER> (defvar *todays-humidity* 70)
*TODAYS-HUMIDITY*
\end{verbatim}

The asterisks are part of the names. They mean nothing to Lisp. It
is a firm convention to name global variables this way so that they
stand out, and you will see why it matters in a moment.

The two differ in one way. If the variable already has a value,
\code{defvar} leaves it alone, and \code{defparameter} replaces it:

\begin{verbatim}
CL-USER> (defparameter *todays-temp* 100)
*TODAYS-TEMP*
CL-USER> (defvar *todays-humidity* 50)
*TODAYS-HUMIDITY*
CL-USER> *todays-temp*
100
CL-USER> *todays-humidity*
70
\end{verbatim}

This matters because you reload files often while you work. Use
\code{defvar} for state that should survive a reload, such as a table
your program has been filling in, and \code{defparameter} for settings
that should follow the source file.

To change the value of a variable, use \code{setf}\index{setf@\texttt{setf}}:

\begin{verbatim}
CL-USER> (setf *todays-humidity* 30)
30
CL-USER> *todays-humidity*
30
\end{verbatim}

You will also see \code{setq}\index{setq@\texttt{setq}} in older code. It does the same thing
for variables. \code{Setf} is more general, as later sections show,
so it is the one to learn.

\subsubsection{Local Variables}

Local variables\index{local variables} are introduced with \code{let}\index{let@\texttt{let}}:

\begin{verbatim}
CL-USER> (let ((a 20)
               (b 30))
           (+ a b))
50
\end{verbatim}

\code{Let} takes a list of bindings and then a body. Each binding is
a pair: a variable name and an expression for its initial value. The
body is one or more expressions, which can use the variables, and the
value of the last one is the value of the whole \code{let}. Outside
the \code{let}, the variables do not exist.

The initial values in a \code{let} are all computed before any of the
variables exists, so one binding cannot refer to another. When you
want that, use \code{let*}\index{let*@\texttt{let*}}, which makes the bindings one after
another:

\begin{verbatim}
CL-USER> (let* ((a 20)
                (b (* a 2)))
           (+ a b))
60
\end{verbatim}

Variables made by \code{let} have \emph{lexical scope}\index{lexical scope}: they can be
seen only by the code written inside the body of the \code{let}. You
can tell where a variable is visible by looking at the text of the
program.

\subsubsection{Dynamic Scope}

Global variables behave differently, and usefully so. If you bind a
global variable with \code{let}, the new value is seen not only by the
code inside the \code{let} but by every function called from there,
for as long as the \code{let} is running. Afterwards the old value is
back.

\begin{verbatim}
CL-USER> (defun report-humidity ()
           (format nil "Humidity is ~a%" *todays-humidity*))
REPORT-HUMIDITY
CL-USER> (report-humidity)
"Humidity is 30%"
CL-USER> (let ((*todays-humidity* 50))
           (report-humidity))
"Humidity is 50%"
CL-USER> (report-humidity)
"Humidity is 30%"
\end{verbatim}

This is called \emph{dynamic scope}\index{dynamic scope}, and global variables are also
known as \emph{special variables}\index{special variables} because of it. It is the clean way
to change a setting for the length of one operation. You do not have
to remember to put the old value back, even if an error interrupts the
work, and in a program with several threads the temporary value is
seen only by the thread that made it. The standard output stream, for
instance, is a special variable, so you can redirect the output of a
whole body of code by binding it.

It is also why the asterisks matter: they tell you at a glance that a
\code{let} is making a dynamic binding and not a lexical one.

\section{The List as a Data Structure}

This section presents the basic operators for working with lists.

\subsection{Accessing the Elements of a List}

The functions \code{first}\index{first@\texttt{first}} through \code{tenth} return the elements
of a list by position, and \code{nth} takes the position as a number,
counting from zero:

\begin{verbatim}
CL-USER> (first '(a b c))
A
CL-USER> (second '(a b c))
B
CL-USER> (nth 0 '(a b c))
A
CL-USER> (nth 2 '(a b c))
C
\end{verbatim}

\subsection{The ``Rest'' of the Story}

A very common pattern is to do something with the first element of a
list and then do the same with what remains. The function
\code{rest}\index{rest@\texttt{rest}} returns the list without its first element:

\begin{verbatim}
CL-USER> (rest '(a b c))
(B C)
\end{verbatim}

In older code, and in a good deal of current code, you will see
\code{first} and \code{rest} written as \code{car}\index{car@\texttt{car}} and \code{cdr}\index{cdr@\texttt{cdr}},
names that survive from the machine Lisp was first written for. They
mean the same thing.

\subsection{The Empty List}

The symbol \code{nil}\index{nil@\texttt{nil}} is the empty list, which can also be written
\code{()}. \code{Nil} evaluates to itself, so it needs no quote:

\begin{verbatim}
CL-USER> nil
NIL
CL-USER> ()
NIL
\end{verbatim}

\code{Nil} is also what Lisp uses for ``false.'' The symbol
\code{t}\index{t@\texttt{t}} is the standard value for ``true,'' but in any test,
\emph{every value other than \code{nil} counts as true}. Many
functions take advantage of this, as you will see.

The function \code{null}\index{null@\texttt{null}} tests for the empty list:

\begin{verbatim}
CL-USER> (null '(a b c))
NIL
CL-USER> (null nil)
T
\end{verbatim}

\subsection{Are You a List?}

\code{Listp} tests whether something is a list. Functions that ask a
yes-or-no question are called \emph{predicates}\index{predicates}, and by convention
their names end in \code{p}:

\begin{verbatim}
CL-USER> (listp '(pontiac cadillac chevrolet))
T
CL-USER> (listp 99)
NIL
CL-USER> (listp nil)
T
\end{verbatim}

The last result is \code{T} because \code{nil} is a list: the empty
one.

\subsection{The Conditional If}

Before going on with lists we need a way to make a choice.
\code{If}\index{if@\texttt{if}} takes three expressions: a test, a ``then'' and an ``else.''
It evaluates the test. If the result is true, it evaluates the
second expression and returns its value; otherwise it evaluates the
third and returns that. Only one of the two is ever evaluated, which
is why \code{if} cannot be a function.

\begin{verbatim}
CL-USER> (if (> 3 4)
             "yes"
             "no")
"no"
CL-USER> (if (listp '("Chicago" "Detroit" "Toronto"))
             "it is a list"
             "it is not a list")
"it is a list"
\end{verbatim}

Notice that \code{if} returns a value. In Lisp there is no division
between statements and expressions: everything returns a value and can
be used inside something else.

\subsection{Length of a List}

\code{Length}\index{length@\texttt{length}} returns the number of elements in a list:

\begin{verbatim}
CL-USER> (length '(gm ford stellantis volkswagen))
4
CL-USER> (length nil)
0
\end{verbatim}

Like most of Common Lisp, \code{length} could be written in Common
Lisp. Here is a simple version:

\begin{verbatim}
(defun our-length (list)
  (if (null list)
      0
      (+ (our-length (rest list)) 1)))
\end{verbatim}

In English: the length of an empty list is zero, and the length of any
other list is one more than the length of its rest. A function that
calls itself like this is \emph{recursive}\index{recursion}, and recursion is a natural
way to express operations on lists.

\subsection{Member of a List}

\code{Member}\index{member@\texttt{member}} tells you whether an item is in a list. If it is not,
\code{member} returns \code{nil}. If it is, \code{member} returns the
part of the list that starts with the item:

\begin{verbatim}
CL-USER> (member 'dallas '(boston san-francisco portland))
NIL
CL-USER> (member 'san-francisco '(boston san-francisco portland))
(SAN-FRANCISCO PORTLAND)
\end{verbatim}

Because anything that is not \code{nil} counts as true, the second
result works perfectly well as a ``yes'' in a test, and it carries
some useful information besides. This is a common idiom in Lisp.

\subsection{Getting Part of a List}

\code{Subseq}\index{subseq@\texttt{subseq}} returns part of a list. It takes the list, the
position to start at, and optionally the position to stop before:

\begin{verbatim}
CL-USER> (subseq '(a b c d) 1 3)
(B C)
CL-USER> (subseq '(a b c d) 1)
(B C D)
\end{verbatim}

\code{Subseq}, \code{length} and many other functions work on any
\emph{sequence}\index{sequences}, which means lists, strings and the vectors of
Section~\ref{sec:arrays}. You saw \code{subseq} applied to a string
at the start of the chapter.

\subsection{Appending Lists}

\code{Append}\index{append@\texttt{append}} takes any number of lists and returns a new list made
of all their elements:

\begin{verbatim}
CL-USER> (defparameter *slides* '(introduction welcome lists functions))
*SLIDES*
CL-USER> (append *slides* '(numbers))
(INTRODUCTION WELCOME LISTS FUNCTIONS NUMBERS)
CL-USER> *slides*
(INTRODUCTION WELCOME LISTS FUNCTIONS)
\end{verbatim}

\code{Append} did not change \code{*slides*}. Like most CL functions,
it returns a new value and leaves its arguments alone. If you want
the variable to change, say so:

\begin{verbatim}
CL-USER> (setf *slides* (append *slides* '(numbers)))
(INTRODUCTION WELCOME LISTS FUNCTIONS NUMBERS)
\end{verbatim}

\subsection{Adding Elements to a List}
\label{sec:cons}

\code{Cons}\index{cons@\texttt{cons}} adds one element to the front of a list:

\begin{verbatim}
CL-USER> (cons 'a '(b c d))
(A B C D)
\end{verbatim}

\code{Cons} is the primitive from which lists are built. A list is a
chain of pairs, called \emph{cons cells}\index{cons cells}, each holding one element
and a pointer to the next pair. \code{Cons} makes one new cell that
points at the existing list, so it is fast, and the old list is not
copied or changed. When you hear that a program ``conses\index{consing} a lot,''
it means that the program allocates a lot of memory, which the garbage
collector will later have to reclaim.

The macros \code{push}\index{push@\texttt{push}} and \code{pop}\index{pop@\texttt{pop}} combine \code{cons} and
\code{rest} with updating a variable, which lets you use a list as a
stack:

\begin{verbatim}
CL-USER> (defparameter *stack* '(b c))
*STACK*
CL-USER> (push 'a *stack*)
(A B C)
CL-USER> (pop *stack*)
A
CL-USER> *stack*
(B C)
\end{verbatim}

\subsection{Removing Elements from a List}

\code{Remove}\index{remove@\texttt{remove}} returns a new list without the occurrences of an item.
\code{Remove-if} and \code{remove-if-not} take a predicate in place of
the item:

\begin{verbatim}
CL-USER> (remove 1000 '(1 2 3 1000 4))
(1 2 3 4)
CL-USER> (remove-if-not #'evenp '(1 2 3 4 5 6))
(2 4 6)
\end{verbatim}

The notation \code{\#\q evenp} means ``the function named
\code{evenp}.'' Section~\ref{sec:functions} explains it.

\subsection{Sorting Lists}

\code{Sort}\index{sort@\texttt{sort}} sorts a sequence, given a function that says whether one
element should come before another:

\begin{verbatim}
CL-USER> (defparameter *numbers* (list 1 3 5 7 2 4 6))
*NUMBERS*
CL-USER> (setf *numbers* (sort *numbers* #'<))
(1 2 3 4 5 6 7)
CL-USER> (setf *numbers* (sort *numbers* #'>))
(7 6 5 4 3 2 1)
\end{verbatim}

\code{Sort} is an exception to the rule that functions leave their
arguments alone. For speed it is allowed to reuse the cells of the
list you give it, so afterwards the original list is in an undefined
state. Always use the value that \code{sort} returns, as above. To
keep the original, sort a copy:

\begin{verbatim}
(defun safe-sort (list predicate)
  (sort (copy-list list) predicate))
\end{verbatim}

Functions that may modify their arguments are called
\emph{destructive}\index{destructive operations}. There are only a few common ones, and their
documentation says so. Never use one on a quoted list written in your
source, because that list is a constant that belongs to the program.

\subsection{Treating a List as a Set}

\code{Union}, \code{intersection} and \code{set-difference} treat two
lists as sets. Sets have no order, so do not rely on the order of the
results:

\begin{verbatim}
CL-USER> (union '(1 2 3) '(2 3 4))
(4 1 2 3)
CL-USER> (intersection '(1 2 3) '(2 3 4))
(3 2)
CL-USER> (set-difference '(1 2 3 4) '(2 3))
(4 1)
\end{verbatim}

\subsection{Mapping a Function to a List}

When you have a list and want another list of the same length, with
each element transformed, use \code{mapcar}\index{mapcar@\texttt{mapcar}}. It takes a function and
a list, calls the function on each element, and returns a list of the
results:

\begin{verbatim}
CL-USER> (defun twice (num)
           (* num 2))
TWICE
CL-USER> (mapcar #'twice '(1 2 3 4))
(2 4 6 8)
\end{verbatim}

Given several lists, \code{mapcar} takes one argument from each:

\begin{verbatim}
CL-USER> (mapcar #'+ '(1 2 3) '(10 20 30))
(11 22 33)
\end{verbatim}

\code{Reduce}\index{reduce@\texttt{reduce}} combines all the elements of a list into one value,
using a function of two arguments:

\begin{verbatim}
CL-USER> (reduce #'+ '(1 2 3 4))
10
\end{verbatim}

\code{Mapcar}, \code{remove-if-not} and \code{reduce} are the map,
filter and reduce of other languages.

\subsection{Property Lists}

A property list\index{property lists}, or \emph{plist}, is a simple way to keep named
values. It is a list in which the odd-numbered items are names,
usually keywords, and each is followed by its value. \code{Getf}\index{getf@\texttt{getf}}
looks up a name:

\begin{verbatim}
CL-USER> (getf '(:michigan "Lansing"
                 :illinois "Springfield"
                 :pennsylvania "Harrisburg")
               :illinois)
"Springfield"
\end{verbatim}

Plists are fine for a handful of items. For large collections, use
the hash tables of Section~\ref{sec:hash-tables}.

\section{Control of Execution}

Controlling execution in CL comes down to controlling which
expressions get evaluated.

\subsection{If, When and Unless}

\code{If} takes exactly one ``then'' expression and one ``else''
expression. To do several things in either branch, group them with
\code{progn}\index{progn@\texttt{progn}}, which evaluates its arguments in order and returns the
value of the last:

\begin{verbatim}
(if (eql status :all-systems-go)
    (progn (broadcast-countdown)
           (launch-rocket))
    (progn (broadcast-apology)
           (shut-down-rocket)))
\end{verbatim}

When there is no ``else,'' use \code{when}\index{when@\texttt{when}}. It takes a test and any
number of expressions, with no need for \code{progn}.
\code{Unless}\index{unless@\texttt{unless}} is the opposite: it evaluates its body if the test is
false.

\begin{verbatim}
CL-USER> (when (> 4 3)
           (print "four is more")
           :done)

"four is more"
:DONE
CL-USER> (unless (> 4 3)
           :never)
NIL
\end{verbatim}

\subsection{Logical Operators}

\code{And}\index{and@\texttt{and}} evaluates its arguments from left to right. It stops and
returns \code{nil} as soon as one of them is false. If none is, it
returns the value of the last one. \code{Or}\index{or@\texttt{or}} stops at the first
argument that is true and returns that value:

\begin{verbatim}
CL-USER> (and (listp '(chicago detroit boston new-york))
              (listp 3.1415))
NIL
CL-USER> (or (listp '(chicago detroit boston new-york))
             (listp 3.1415))
T
CL-USER> (or nil 42 99)
42
\end{verbatim}

The last example shows a common way to supply a default:
\code{(or value default)} returns \code{value} unless it is
\code{nil}.

\code{Not}\index{not@\texttt{not}} returns \code{t} if its argument is \code{nil}, and
\code{nil} otherwise.

\subsection{Cond}

For a chain of tests, where another language would say ``if \ldots{}
else if \ldots{} else,'' use \code{cond}\index{cond@\texttt{cond}}. Each clause is a list holding
a test followed by expressions. \code{Cond} tries the tests in order,
and at the first true one it evaluates that clause's expressions and
returns. A final clause whose test is \code{t} catches everything
else:

\begin{verbatim}
CL-USER> (let ((n 4))
           (cond ((> n 10) "It's a lot")
                 ((= n 10) "It's kind of a lot")
                 (t "It's not a lot")))
"It's not a lot"
\end{verbatim}

With \code{cond} we can write our own version of \code{member}, very
close to how you would describe it in English:

\begin{verbatim}
(defun our-member (elem list)
  (cond ((null list) nil)
        ((eql elem (first list)) list)
        (t (our-member elem (rest list)))))
\end{verbatim}

``If the list is empty, the answer is no. If the item is the same as
the first element, return the list. Otherwise look in the rest.''

\subsection{Case}

To choose among fixed values, \code{case}\index{case@\texttt{case}} is more concise than
\code{cond}:

\begin{verbatim}
CL-USER> (let ((color :red))
           (case color
             (:blue "Blue is okay")
             (:red "Red is actually her favorite color")
             (otherwise "We have never heard of that!")))
"Red is actually her favorite color"
\end{verbatim}

\code{Case} compares with \code{eql}, so it works with symbols,
keywords, integers and characters, and not with strings.
Section~\ref{sec:equality} says why. \code{Ecase}\index{ecase@\texttt{ecase}} is like
\code{case} but signals an error if nothing matches, which is what
you want when the list of cases is supposed to be complete.

\subsection{Iteration}

Lisp has recursion and mapping, and it also has ordinary loops\index{iteration}.
\code{Dolist}\index{dolist@\texttt{dolist}} runs its body once for each element of a list, and
\code{dotimes}\index{dotimes@\texttt{dotimes}} runs it a given number of times, counting from zero:

\begin{verbatim}
CL-USER> (dolist (city '("Detroit" "Delft" "Tokyo"))
           (format t "Hello, ~a!~%" city))
Hello, Detroit!
Hello, Delft!
Hello, Tokyo!
NIL
CL-USER> (let ((result 0))
           (dotimes (n 10)
             (setf result (+ result n)))
           result)
45
\end{verbatim}

The \code{loop}\index{loop@\texttt{loop}} macro is a small language of its own for iteration.
It reads almost like English, and it can collect, sum, count and
filter as it goes:

\begin{verbatim}
CL-USER> (loop for n from 1 to 5 collect (* n n))
(1 4 9 16 25)
CL-USER> (loop for elem in '(4 5 6 7) sum elem)
22
CL-USER> (loop for n from 1 to 10
               when (evenp n) collect n)
(2 4 6 8 10)
\end{verbatim}

\code{Loop} can do a great deal more. The references in Chapter~5
describe it in full. To leave any of these loops early, call
\code{return}.

\section{Functions as Objects}
\label{sec:functions}

\subsection{Named Functions}

A function in CL is an object\index{functions!as objects}, like a number or a list. It can be
stored in a variable, passed to another function, and returned as a
result. To get the function object that a name refers to, write
\code{\#\q} before the name:

\begin{verbatim}
CL-USER> #'+
#<FUNCTION +>
\end{verbatim}

A function object has no printed form that could be read back in, so
Lisp prints a description between \code{\#<} and \code{>}. You will
see this notation for many kinds of object.

To call a function object, use \code{funcall}\index{funcall@\texttt{funcall}}, giving the arguments
one by one, or \code{apply}\index{apply@\texttt{apply}}, giving them as a list:

\begin{verbatim}
CL-USER> (funcall #'+ 1 2)
3
CL-USER> (apply #'+ '(1 2 3))
6
\end{verbatim}

\subsection{Functional Arguments}

You have already passed functions to \code{mapcar}, \code{sort},
\code{reduce} and \code{remove-if-not}. Functions that take other
functions as arguments are called higher-order functions\index{higher-order functions}, and CL uses
them everywhere. Your own functions qualify as soon as you define
them:

\begin{verbatim}
CL-USER> (defun add-3 (num)
           (+ num 3))
ADD-3
CL-USER> (mapcar #'add-3 '(2 3 4))
(5 6 7)
\end{verbatim}

\subsection{Anonymous Functions}

Often a function is needed in just one place, and it is not worth
giving it a name. \code{Lambda}\index{lambda@\texttt{lambda}} makes a function on the spot. It
looks like \code{defun} without the name:

\begin{verbatim}
CL-USER> (mapcar (lambda (num) (+ num 3))
                 '(2 3 4))
(5 6 7)
CL-USER> (sort (list '(4 "Buffy") '(2 "Keiko") '(1 "Judy") '(3 "Aruna"))
               (lambda (x y)
                 (< (first x) (first y))))
((1 "Judy") (2 "Keiko") (3 "Aruna") (4 "Buffy"))
\end{verbatim}

A function made by \code{lambda} can use the variables that are
visible where it is written, and it keeps them for as long as it
lives. Such a function is called a \emph{closure}\index{closures}:

\begin{verbatim}
CL-USER> (defun make-adder (n)
           (lambda (x) (+ x n)))
MAKE-ADDER
CL-USER> (defparameter *add-10* (make-adder 10))
*ADD-10*
CL-USER> (funcall *add-10* 5)
15
CL-USER> (mapcar (make-adder 100) '(1 2 3))
(101 102 103)
\end{verbatim}

Each call to \code{make-adder} returns a new function that remembers
its own \code{n}.

\subsection{Optional Arguments}

Parameters after the symbol \code{\&optional}\index{optional arguments} in a parameter list
may be left out by the caller. Each can be given a default value;
without one, the default is \code{nil}:

\begin{verbatim}
CL-USER> (defun greeting-list (&optional (username "Jake"))
           (list 'hello 'there username))
GREETING-LIST
CL-USER> (greeting-list)
(HELLO THERE "Jake")
CL-USER> (greeting-list "Joe")
(HELLO THERE "Joe")
\end{verbatim}

\subsection{Keyword Arguments}

Parameters after \code{\&key}\index{keyword arguments} are optional too, and the caller passes
them by name, in any order. They are written the same way:

\begin{verbatim}
(defun greeting-list (&key (username "Jake")
                           (greeting "how are you?"))
  (list 'hello 'there username greeting))
\end{verbatim}

\noindent
and called by putting a keyword before each value:

\begin{verbatim}
CL-USER> (greeting-list :greeting "how have you been?")
(HELLO THERE "Jake" "how have you been?")
CL-USER> (greeting-list :greeting "how have you been?" :username "Joe")
(HELLO THERE "Joe" "how have you been?")
\end{verbatim}

Keyword arguments make calls easy to read, and they let you add an
option to a function later without disturbing its existing callers.
The built-in functions use them heavily. \code{Sort}, for instance,
takes a \code{:key} function that picks out the part of each element
to compare, so the sort above could have been written
\code{(sort people \#\q< :key \#\q first)}.

A parameter after \code{\&rest}\index{rest arguments} receives all the remaining arguments
as a list, which is how you write a function that takes any number of
them:

\begin{verbatim}
CL-USER> (defun average (&rest numbers)
           (/ (reduce #'+ numbers) (length numbers)))
AVERAGE
CL-USER> (average 1 2 3 4)
5/2
\end{verbatim}

\subsection{Multiple Values}

A function can return more than one value\index{multiple values}. \code{Floor}, for
example, divides and returns both the quotient and the remainder:

\begin{verbatim}
CL-USER> (floor 7 2)
3
1
\end{verbatim}

If you use the call in the ordinary way, you get the first value and
the others are quietly dropped, so \code{(+ (floor 7 2) 10)} is 13.
To receive them all, use \code{multiple-value-bind}\index{multiple-value-bind@\texttt{multiple-value-bind}}:

\begin{verbatim}
CL-USER> (multiple-value-bind (quotient remainder)
             (floor 7 2)
           (list quotient remainder))
(3 1)
\end{verbatim}

Your own functions return several values by calling \code{values}:
\code{(values 3 1)}.

\section{Input, Output, Streams, and Strings}

Input and output in CL go through \emph{streams}\index{streams}. A stream is a
source or destination of data: the terminal, a file, a network
connection, or a string in memory. Two special variables hold the
defaults: \code{*standard-input*} and \code{*standard-output*}.

\subsection{Format}

\code{Format}\index{format@\texttt{format}} is the function you will use for most output. It takes
a destination, a control string, and arguments that are filled in
where the control string has directives. Directives start with a
tilde:

\begin{verbatim}
CL-USER> (format t "Hello there, ~a!~%" 'bob)
Hello there, BOB!
NIL
\end{verbatim}

The destination \code{t} means standard output. The directive
\code{\td a} prints the next argument in a form meant for people, and
\code{\td\%} starts a new line. \code{Format} printed its text and
then returned \code{nil}, which the REPL showed.

If the destination is \code{nil}, \code{format} prints nothing and
returns the text as a string. This is the usual way to build a string
from pieces:

\begin{verbatim}
CL-USER> (format nil "Hello there, ~a!" 'bob)
"Hello there, BOB!"
\end{verbatim}

A few more directives will cover most of what you need:
\code{\td s} prints an argument the way you would type it in a
program, \code{\td d} prints an integer, \code{\td ,2f} prints a
number with two decimal places, and \code{\td\{} \ldots{}
\code{\td\}} repeats over a list:

\begin{verbatim}
CL-USER> (format nil "~a and ~s" "a string" "a string")
"a string and \"a string\""
CL-USER> (format nil "~d items at ~,2f each" 3 4.5)
"3 items at 4.50 each"
CL-USER> (format nil "~{~a~^, ~}" '(1 2 3))
"1, 2, 3"
\end{verbatim}

There are many others. One prints numbers in English and another in
Roman numerals:

\begin{verbatim}
CL-USER> (format nil "~r" 2026)
"two thousand twenty-six"
CL-USER> (format nil "~@r" 2026)
"MMXXVI"
\end{verbatim}

\subsection{Print and Princ}

\code{Print}\index{print@\texttt{print}} outputs one object in the form that Lisp could read back
in, on a new line, and returns the object. \code{Princ}\index{princ@\texttt{princ}} outputs it in
the form meant for people. The difference shows with strings:

\begin{verbatim}
CL-USER> (print "hello")

"hello"
"hello"
CL-USER> (princ "hello")
hello
"hello"
\end{verbatim}

In each case the first thing you see is the output and the second is
the return value, printed by the REPL. \code{Print} returns its
argument, so you can wrap it around any expression in the middle of a
program to see a value go by without changing what the program does.

\subsection{Read}

\code{Read}\index{read@\texttt{read}} is the other half of the REPL: it takes text from a stream
and returns the Lisp object that the text describes. It does not
evaluate anything. The whole reader is available to your programs,
which makes Lisp's own notation a ready-made data format.

\begin{verbatim}
CL-USER> (read-from-string "(a b (c 42) \"text\")")
(A B (C 42) "text")
19
\end{verbatim}

\noindent
(The second value is the position where reading stopped.) Only read
data that you trust this way. For text from outside, read lines with
\code{read-line}\index{read-line@\texttt{read-line}}, which returns one line as a string, or use a
library for a format such as JSON, as in Chapter~4.

\subsection{File Input and Output}

To use a file you open a stream to it. The macro
\code{with-open-file}\index{with-open-file@\texttt{with-open-file}} opens the file, gives you the stream in a
variable for the length of its body, and closes the file when the body
is finished, even if it is left because of an error.

To write a file:

\begin{verbatim}
CL-USER> (with-open-file (out "/tmp/readme.txt"
                              :direction :output
                              :if-exists :supersede)
           (format out "Hello there~%")
           (format out "Second line~%"))
NIL
\end{verbatim}

\noindent
and to read it back, a line at a time:

\begin{verbatim}
CL-USER> (with-open-file (in "/tmp/readme.txt")
           (loop for line = (read-line in nil)
                 while line
                 collect line))
("Hello there" "Second line")
\end{verbatim}

Opening for input is the default, so the second form needs no options.
The \code{nil} given to \code{read-line} tells it to return
\code{nil} at the end of the file, where it would otherwise signal
an error.

\code{Probe-file}\index{probe-file@\texttt{probe-file}} tests whether a file exists.

\subsection{Pathnames}

You can name files with ordinary strings, as above. Lisp turns them
into \emph{pathname}\index{pathnames} objects, which you can take apart and combine
without cutting strings up by hand:

\begin{verbatim}
CL-USER> (defparameter *my-path* (pathname "/tmp/readme.txt"))
*MY-PATH*
CL-USER> *my-path*
#P"/tmp/readme.txt"
CL-USER> (pathname-name *my-path*)
"readme"
CL-USER> (pathname-type *my-path*)
"txt"
CL-USER> (merge-pathnames "notes.txt" *my-path*)
#P"/tmp/notes.txt"
\end{verbatim}

\section{Hash Tables, Arrays, Structures, and Classes}

Lists are convenient, and they are not the answer to everything. CL
has the other data structures you would expect.

\subsection{Hash Tables}
\label{sec:hash-tables}

A hash table\index{hash tables} maps keys to values, like a dictionary or map in other
languages. Lookup stays fast however large the table grows. Make one
with \code{make-hash-table}, look things up with \code{gethash}\index{gethash@\texttt{gethash}}, and
store them with \code{setf} of \code{gethash}:

\begin{verbatim}
CL-USER> (defparameter *ratings* (make-hash-table))
*RATINGS*
CL-USER> (setf (gethash :lisp *ratings*) :excellent)
:EXCELLENT
CL-USER> (gethash :lisp *ratings*)
:EXCELLENT
T
CL-USER> (gethash :cobol *ratings*)
NIL
NIL
\end{verbatim}

This is \code{setf}\index{setf@\texttt{setf}} doing something \code{setq} cannot. Its first
argument is not just a variable but a \emph{place}: an expression
that says where a value lives. To store into anything, you write the
expression that would fetch the value and wrap it in \code{setf}. The
same rule works for hash tables, arrays, objects and parts of lists,
so there is no separate ``put'' function to learn for each.

\code{Gethash} returns two values. The second says whether the key
was found, which lets you tell a missing key from a key whose value is
\code{nil}.

A table made with no options compares keys with \code{eql}, which
suits symbols and numbers. For string keys, ask for \code{equal}:

\begin{verbatim}
CL-USER> (defparameter *capitals* (make-hash-table :test #'equal))
*CAPITALS*
CL-USER> (setf (gethash "Michigan" *capitals*) "Lansing")
"Lansing"
CL-USER> (gethash "Michigan" *capitals*)
"Lansing"
T
\end{verbatim}

\code{Maphash} calls a function on each key and value, and
\code{hash-table-count} returns the number of entries.

\subsection{Arrays}
\label{sec:arrays}

Arrays\index{arrays} hold values in a grid of one or more dimensions, indexed by
integers from zero, with fast access to any element.
\code{Make-array} creates one and \code{aref}\index{aref@\texttt{aref}} reads an element. As
with hash tables, \code{setf} of the same expression stores one:

\begin{verbatim}
CL-USER> (defparameter *grid* (make-array '(3 3) :initial-element 0))
*GRID*
CL-USER> (setf (aref *grid* 0 0) "Dave")
"Dave"
CL-USER> (setf (aref *grid* 0 1) "John")
"John"
CL-USER> *grid*
#2A(("Dave" "John" 0) (0 0 0) (0 0 0))
\end{verbatim}

A one-dimensional array is called a \emph{vector}\index{vectors}. Vectors can be
written directly with \code{\#(} \ldots{} \code{)}, and they can be made
to grow:

\begin{verbatim}
CL-USER> (aref #(10 20 30) 1)
20
CL-USER> (defparameter *v* (make-array 0 :adjustable t :fill-pointer t))
*V*
CL-USER> (vector-push-extend :a *v*)
0
CL-USER> (vector-push-extend :b *v*)
1
CL-USER> *v*
#(:A :B)
\end{verbatim}

Strings are vectors of characters. Vectors and lists together are the
sequences, so \code{length}, \code{subseq}, \code{sort},
\code{remove}, \code{find}, \code{position} and \code{reduce} work on
all of them.

\subsection{Structures}

A structure\index{structures} is a record with named slots. One \code{defstruct}\index{defstruct@\texttt{defstruct}}
form defines the type together with a function to make instances, a
function to read each slot, and a predicate:

\begin{verbatim}
CL-USER> (defstruct point
           x
           y)
POINT
CL-USER> (defparameter *p* (make-point :x 3 :y 4))
*P*
CL-USER> (point-x *p*)
3
CL-USER> (setf (point-y *p*) 10)
10
CL-USER> *p*
#S(POINT :X 3 :Y 10)
\end{verbatim}

Structures are simple and fast. When you need more than that, use
classes.

\subsection{Classes and Methods}

The Common Lisp Object System\index{CLOS}, CLOS, is the object-oriented part of
the language. A class\index{classes} describes a kind of object: its slots, and
the classes it inherits from. \code{Defclass}\index{defclass@\texttt{defclass}} defines one:

\begin{verbatim}
(defclass box ()
  ((length :initarg :length :initform 10 :accessor box-length)
   (width  :initarg :width  :initform 10 :accessor box-width)
   (height :initarg :height :initform 10 :accessor box-height)))
\end{verbatim}

After the name comes the list of superclasses, empty here, and then
the slots. For each slot, \code{:initarg} gives the keyword used to
set it when an object is made, \code{:initform} gives its default
value, and \code{:accessor} names a function for reading the slot,
and for writing it with \code{setf}.

Behavior is defined with \emph{generic functions}\index{generic functions} and
\emph{methods}\index{methods}. A generic function is called like any other
function. What it does depends on the classes of its arguments, and
each method supplies the behavior for one case:

\begin{verbatim}
(defgeneric volume (thing)
  (:documentation "Return the volume of THING."))

(defmethod volume ((self box))
  (* (box-length self)
     (box-width self)
     (box-height self)))
\end{verbatim}

\noindent
\code{Make-instance} creates an object:

\begin{verbatim}
CL-USER> (defparameter *box* (make-instance 'box :height 2))
*BOX*
CL-USER> (volume *box*)
200
CL-USER> (setf (box-width *box*) 5)
5
CL-USER> (volume *box*)
100
\end{verbatim}

A subclass inherits slots and methods:

\begin{verbatim}
CL-USER> (defclass labeled-box (box)
           ((label :initarg :label :accessor box-label)))
#<STANDARD-CLASS COMMON-LISP-USER::LABELED-BOX>
CL-USER> (volume (make-instance 'labeled-box :label "Fragile"))
1000
\end{verbatim}

If you know Java, C++ or Python, notice what is different here.
Methods do not live inside classes. A call is written
\code{(volume box)}, like any function call, and not
\code{box.volume()}. One consequence is that you can add a method for
a class you did not write, including the built-in ones, without
touching it. Another is that a method can be chosen by the classes of
\emph{all} its arguments, not only the first:

\begin{verbatim}
CL-USER> (defgeneric add (a b))
#<STANDARD-GENERIC-FUNCTION COMMON-LISP-USER::ADD (0)>
CL-USER> (defmethod add ((a string) (b string))
           (concatenate 'string a b))
#<STANDARD-METHOD COMMON-LISP-USER::ADD (STRING STRING) {1201ACA633}>
CL-USER> (defmethod add ((a number) (b number))
           (+ a b))
#<STANDARD-METHOD COMMON-LISP-USER::ADD (NUMBER NUMBER) {1201B1BD13}>
CL-USER> (add "hello " "there")
"hello there"
CL-USER> (add 3 4)
7
\end{verbatim}

CLOS goes a good deal further than this, with multiple inheritance,
methods that run before, after and around others, and classes that can
be redefined while instances of them exist. The existing objects are
updated to match. It fits the way Lisp is developed, and it is one of
the reasons a Lisp program can stay up while it is being changed.

\section{Macros}
\label{sec:macros}

You have been using macros\index{macros} since the first \code{defun}:
\code{when}, \code{cond}, \code{dolist}, \code{loop}, \code{setf} and
\code{with-open-file} are all macros. This section shows what they
are and how to write a small one. Writing macros well is a subject
for the longer books. Reading them is something you need from the
first day.

A function receives the \emph{values} of its arguments. A macro
receives its arguments as unevaluated \emph{code}, as the lists and
symbols you typed, and returns a new piece of code to be used in their
place. It is a function from code to code, run by the compiler.
Because Lisp code is made of lists, a macro is written with the same
list operations as any other program.

Suppose \code{unless} did not exist. A function could not do its job,
because a function's arguments are all evaluated before it is called,
and the point of \code{unless} is that the body sometimes must not be.
A macro can:

\begin{verbatim}
(defmacro our-unless (test &body body)
  `(if ,test
       nil
       (progn ,@body)))
\end{verbatim}

\code{Defmacro}\index{defmacro@\texttt{defmacro}} looks like \code{defun}. The parameter \code{test}
receives the first expression in the call, and \code{\&body} collects
the rest as a list. The body builds the replacement code with the
backquote\index{backquote} notation: a backquote is like a quote, except that inside
it a comma means ``put the value of this here,'' and comma-at means
``splice the elements of this list in here.'' The result is a
template with holes.

You can ask Lisp what a macro call turns into, with
\code{macroexpand-1}\index{macroexpand-1@\texttt{macroexpand-1}}:

\begin{verbatim}
CL-USER> (macroexpand-1 '(our-unless (> 3 4) (print "no") :done))
(IF (> 3 4)
    NIL
    (PROGN (PRINT "no") :DONE))
T
\end{verbatim}

\noindent
and the new operator works like a built-in one:

\begin{verbatim}
CL-USER> (our-unless (> 3 4)
           (print "three is not more")
           :done)

"three is not more"
:DONE
\end{verbatim}

The built-in macros are no different. In SBCL:

\begin{verbatim}
CL-USER> (macroexpand-1 '(when (> 4 3) :done))
(IF (> 4 3)
    :DONE)
T
\end{verbatim}

When you meet a macro you do not know, expand it. Your editor has a
key for this (Table~\ref{tab:slime}).

This is what Chapter~1 meant by extending the language from the
inside. The built-in \code{with-open-file} is a macro that someone
wrote, and in the same way you can write your own, such as a
\code{with-connection} or a \code{define-report}, for the needs of
your program. A good rule for beginners is to
write a function whenever a function will do, and to reach for a macro
when you need to control evaluation or to remove a pattern that keeps
repeating in your code.

\section{Conditions and Error Handling}
\label{sec:conditions}

Chapter~2 showed what happens at the REPL when something goes wrong:
you land in the debugger. A finished program usually has to deal with
problems itself. In CL an error is represented by an object called a
\emph{condition}\index{conditions}, and conditions are organized into classes, such as
\code{division-by-zero}, \code{file-error} and \code{type-error}, all
of which are kinds of \code{error}.

\code{Handler-case}\index{handler-case@\texttt{handler-case}} runs an expression and says what to do if a
condition of a given class is signaled inside it. It corresponds to
try and catch in other languages:

\begin{verbatim}
CL-USER> (handler-case (/ 10 0)
           (division-by-zero ()
             :infinity))
:INFINITY
\end{verbatim}

To signal an error of your own, call \code{error}\index{error@\texttt{error}} with a message in
the style of \code{format}:

\begin{verbatim}
(defun parse-age (string)
  (let ((age (parse-integer string)))
    (when (minusp age)
      (error "An age cannot be negative: ~a" age))
    age))
\end{verbatim}

\begin{verbatim}
CL-USER> (parse-age "42")
42
CL-USER> (handler-case (parse-age "forty")
           (error (condition)
             (format t "Could not read the age: ~a~%" condition)
             nil))
Could not read the age: junk in string "forty"
NIL
\end{verbatim}

The handler names the class \code{error}, so it catches both the
error that \code{parse-integer} signaled here and the one
\code{parse-age} would signal for a negative number. The variable
\code{condition} holds the condition object, and printing it gives
its message.

To make sure something happens on the way out, whether or not there
was an error, use \code{unwind-protect}\index{unwind-protect@\texttt{unwind-protect}}. Its first expression is the
work and the rest are the cleanup:

\begin{verbatim}
CL-USER> (unwind-protect
              (progn (print "working") :result)
           (print "cleaning up"))

"working"
"cleaning up"
:RESULT
\end{verbatim}

\noindent
This is what \code{with-open-file} uses to close the file.

The condition system goes further than try and catch in one important
way. In most languages, by the time an error is caught, the code that
ran into it has been abandoned. In CL the handler runs first, while
that code is still waiting, and it can choose one of the restarts\index{restarts} that
the code has offered: use this value in place of the bad one, skip
this record, try again. The low-level code knows the possible ways to
recover and the high-level code decides which one to use. The
restarts you saw in the debugger in Section~\ref{sec:debugging} are
the same mechanism with a person making the choice. The operators to
look up when you need this are \code{handler-bind} and
\code{restart-case}.

\section{Packages}
\label{sec:packages}

A package\index{packages} is a namespace for symbols. Packages keep the names in one
library from colliding with the same names in another, and they let a
library say which of its names are meant for others to use.

There is always a current package, held in the special variable
\code{*package*}, and the REPL shows it in the prompt. A new Lisp
starts in the package \code{COMMON-LISP-USER}, or \code{CL-USER} for
short, which is a scratch area. A program of any size should have a
package of its own. \code{Defpackage}\index{defpackage@\texttt{defpackage}} defines one:

\begin{verbatim}
(defpackage :inventory
  (:use :common-lisp)
  (:export #:add-item #:*items*))
\end{verbatim}

This says that code in the \code{inventory} package can use all the
standard names of the language, which live in the package
\code{COMMON-LISP}, and that two of its own names are for the outside
world. (The \code{\#:} prefix just writes a name without creating a
symbol in the current package. You will also see keywords and strings
used here.)

Each source file then begins by saying which package it is in:

\begin{verbatim}
(in-package :inventory)

(defvar *items* nil)

(defun add-item (item)
  (push item *items*))

(defun helper ()
  :internal)
\end{verbatim}

From another package, an exported name is written with the package
name and one colon:

\begin{verbatim}
CL-USER> (inventory:add-item :widget)
(:WIDGET)
CL-USER> inventory:*items*
(:WIDGET)
\end{verbatim}

A name that was not exported can still be reached by using two
colons, as in \code{inventory::helper}. The second colon is a warning to
yourself: this name is not part of the library's interface and may
change. It is handy for poking around at the REPL and has no place in
finished code.

If you are in a different package from the one you think you are in,
Lisp will tell you that functions you have just defined do not exist.
When that happens, look at the prompt.

Package names should be distinctive, since all packages share one list
of names, and long names are tedious to type. A package can give
another package a short local nickname\index{package-local nicknames} for its own use:

\begin{verbatim}
(defpackage :report
  (:use :common-lisp)
  (:local-nicknames (:a :alexandria)))
\end{verbatim}

\noindent
after which code in \code{report} can write \code{a:flatten}. Local
nicknames came after the standard, and every implementation mentioned
in this book supports them.

\subsection{The Keyword Package}

Keywords\index{keywords}, the symbols written with a leading colon, live in a package
of their own named \code{KEYWORD}. That is why \code{:red} means the
same thing wherever it is written, whatever the current package, and
why keywords are used for option names and for labels in data.

\section{Common Stumbling Blocks}

This section goes over the things that most often trip up newcomers.

\subsection{Quotes}

A symbol without a quote is a variable. With a quote it is the symbol
itself:

\begin{verbatim}
CL-USER> (defparameter *x* 1)
*X*
CL-USER> *x*
1
CL-USER> '*x*
*X*
\end{verbatim}

A list without a quote is a call. With a quote it is the list:

\begin{verbatim}
CL-USER> (first '(+ 1 2))
+
CL-USER> (first (+ 1 2))
; Evaluation aborted on #<SIMPLE-TYPE-ERROR expected-type: LIST datum: 3>.
\end{verbatim}

In the second, \code{(+ 1 2)} was evaluated first, and \code{first}
was handed the number 3.

Numbers, strings and keywords evaluate to themselves and never need a
quote.

\subsection{Function Argument Lists}

When you define a function, the parameters go in a list of their own.
When you call it, the arguments follow the name directly:

\begin{verbatim}
CL-USER> (defun square (num)
           (* num num))
SQUARE
CL-USER> (square 4)
16
\end{verbatim}

If you are used to writing \code{square(4)}, sooner or later you will
type \code{(square (4))}. Lisp will try to call a function named
\code{4} and complain of an illegal function call.

\subsection{Symbols vs. Strings}

The symbol \code{AAA} and the string \code{"AAA"} are different kinds
of thing. A symbol is a name: there is exactly one symbol with a
given name in a given package, and comparing two symbols is a single
machine instruction. A string is a piece of text: two strings with
the same characters are still two objects.

\begin{verbatim}
CL-USER> (eql 'aaa 'aaa)
T
CL-USER> (eql "bbb" "bbb")
NIL
CL-USER> (string= "bbb" "bbb")
T
\end{verbatim}

Use symbols and keywords for things your program needs to tell apart,
such as the states of a machine or the options of a function. Use
strings for text that people read or that comes from outside. The
string operations do not apply to symbols:

\begin{verbatim}
CL-USER> (defparameter *sentence* "Jane dates only Lisp programmers")
*SENTENCE*
CL-USER> (char *sentence* 3)
#\e
CL-USER> (subseq *sentence* 0 4)
"Jane"
CL-USER> (position #\Space *sentence*)
4
CL-USER> (subseq *sentence* 11)
"only Lisp programmers"
\end{verbatim}

Remember also that the reader converts symbol names to upper case, so
the name of the symbol you type as \code{hello} is the string
\code{"HELLO"}.

\subsection{Equality}
\label{sec:equality}

CL has several functions for testing equality\index{equality}, because ``the same''
means different things for different kinds of object. A rule of
thumb:

\begin{itemize}
\item \code{eql}\index{eql@\texttt{eql}} asks whether two things are the very same object.
Use it for symbols and keywords. It is also true for two integers of
the same value and for two identical characters. It is the default
test used by \code{member}, \code{case}, \code{remove}, \code{find}
and hash tables.

\item \code{=} compares numbers by value, whatever their types.

\item \code{equal}\index{equal@\texttt{equal}} compares structure: two lists with equal elements
are \code{equal}, and so are two strings with the same characters.

\item \code{string=} compares strings, and \code{string-equal} does so
ignoring case.
\end{itemize}

\begin{verbatim}
CL-USER> (eql 3 3.0)
NIL
CL-USER> (= 3 3.0)
T
CL-USER> (eql '(a b c) (list 'a 'b 'c))
NIL
CL-USER> (equal '(a b c) (list 'a 'b 'c))
T
CL-USER> (string-equal "bbb" "BBB")
T
\end{verbatim}

The functions that use \code{eql} by default take a \code{:test}
argument to change it. Forgetting this with strings is the classic
mistake:

\begin{verbatim}
CL-USER> (member "blue" '("red" "blue" "green" "yellow"))
NIL
CL-USER> (member "blue" '("red" "blue" "green" "yellow")
                 :test #'string=)
("blue" "green" "yellow")
\end{verbatim}

Floating-point numbers that ought to be equal often differ by a tiny
amount, in Lisp as in every language. Compare them with a tolerance:

\begin{verbatim}
CL-USER> (= (+ 0.1d0 0.2d0) 0.3d0)
NIL
CL-USER> (defun nearly-equal-p (a b &optional (tolerance 1d-9))
           (< (abs (- a b)) tolerance))
NEARLY-EQUAL-P
CL-USER> (nearly-equal-p (+ 0.1d0 0.2d0) 0.3d0)
T
\end{verbatim}

\subsection{Distinguishing Macros from Functions}

Macros and functions are called with the same syntax, and it rarely
matters which is which. Assume a function unless you have a reason to
think otherwise. Most macros are easy to spot:

\begin{itemize}
\item the familiar control operators, such as \code{when},
\code{cond}, \code{case}, \code{and}, \code{or} and \code{loop};
\item names that begin with \code{def}, such as \code{defun} and
\code{defparameter};
\item names that begin with \code{with-}, such as
\code{with-open-file};
\item names that begin with \code{do}, such as \code{dolist} and
\code{dotimes}.
\end{itemize}

A small number of operators, including \code{if}, \code{let},
\code{quote}, \code{setq} and \code{progn}, are neither. They are the
special operators that the language is built on, and they behave like
macros as far as you are concerned.

If you are in doubt, ask the running Lisp:

\begin{verbatim}
CL-USER> (macro-function 'list)
NIL
CL-USER> (special-operator-p 'if)
T
\end{verbatim}

\noindent
\code{Macro-function} returns \code{nil} for anything that is not a
macro. \code{Describe} will tell you all of this at once, and your
editor's key for it is in Table~\ref{tab:slime}.

\subsection{Modifying What You Should Not}

A quoted list in your source is a constant. Lisp is allowed to store
one copy of it and share it, so changing it with a destructive
operation such as \code{sort} can change your program in ways that are
very hard to track down. If you are going to modify a list, build it
with \code{list} or \code{copy-list}, as the \code{sort} examples in
this chapter did.

Destructive operations exist for speed, and you seldom need them. A
program that builds new lists and lets the garbage collector take the
old ones is easier to get right, and usually fast enough. Write it
that way first and measure with \code{time} before you change
anything.


% Chapter 4
\chapter{Interfaces}

CL provides many mechanisms for interfacing with the ``outside world.'' This chapter will
introduce some of the more important and practical ones available as part of Allegro CL\index{Allegro CL},
or as extensions to Allegro CL through third parties. These techniques include:

\begin{itemize}
\item Direct operating system calls
\item Foreign Function Interface (FFI) for calling functions written in other languages
\item CORBA\index{CORBA}
\item TCP/IP Sockets
\item Database interfaces
\item Web Serving
\end{itemize}

\section{The Operating System}

\subsection{System Calls}

CL provides the function \texttt{excl:run-shell-command} for making direct calls to the operating
system. Note that the \texttt{excl:} package qualifier is used, meaning that this call is an extension
to Common Lisp\index{Common Lisp} and is specific to Allegro CL. Other CL platforms provide similar functionality. Calling OS commands directly is generally a convenient way of interfacing with the
Unix (Linux) operating system or with legacy applications. It is also a quick way to accomplish tasks such as file copying, moving, and deletion without using CL equivalents such
as \texttt{rename-file} and \texttt{delete-file}.

\texttt{Run-shell-command} takes a string representing the desired command, and any optional
keyword arguments as specified in the Allegro CL documentation.

\begin{verbatim}
CL-USER(2): (excl:run-shell-command "ls -l")
total 24
-rwxr-xr-x  1 dcooper users 2863 May  9 14:06 compile-system.cl
-rwxr-xr-x  1 dcooper users 1856 May  8 18:13 example.cl
-rw-r--r--  1 dcooper users  269 May  8 17:48 load-system.cl
-rw-r--r--  1 dcooper users 2152 May  8 18:10 readme.txt
-rwxr-xr-x  1 dcooper users  356 May  8 18:10 start.sh
0
CL-USER(3): (excl:run-shell-command "cp readme.txt readme.bak")
0
\end{verbatim}

\subsection{Environment Variables}

To access values from Unix environment variables, you can use the function \texttt{sys:getenv},
which takes a string representing the name of the environment variable, and returns a string
representing the current value (or NIL if the variable does not exist):

\begin{verbatim}
CL-USER(4): (sys:getenv "SHELL")
"/bin/bash"
\end{verbatim}

You can also set the value of environment variables using \texttt{setf}\index{setf@\texttt{setf}}:

\begin{verbatim}
CL-USER(5): (setf (sys:getenv "MYVAR") "yes")
"yes"
CL-USER(6): (sys:getenv "MYVAR")
"yes"
\end{verbatim}

\section{Foreign Function Interface}

CL provides the capability to call functions written in other languages, most notably the C
language. Each CL platform provides its own Foreign Function Interface (``FFI''), and the
details of this are unfortunately not standardized between platforms. To learn the details
of FFI for Allegro CL, you should consult Allegro CL's documentation.

In general, the steps involved in calling foreign functions from CL are:

\begin{enumerate}
\item Define the Lisp names and types for the C functions and any supporting data structures

\item Load the shared library (\texttt{.so} file on Unix/Linux, or \texttt{.dll} on Windows) which contains
the compiled object code for the functions

\item Call the functions
\end{enumerate}

The FFI capability is extremely useful for ``wrapping'' existing libraries of C code, and/or
for giving CL applications convenient access to low-level system calls. The Franz Inc. website
contains documentation for a large library of Unix system calls which are pre-wrapped and
available from CL.

\section{CORBA}

CORBA\footnote{CORBA stands for Common Object Request Broker Architecture} is a standard for writing distributed object-oriented applications which
can communicate over a network. A CORBA-wrapped application can be written in any
language which supports CORBA. These applications can then communicate with each
other, regardless of the language in which they are written.

CorbaLink is a CORBA implementation for Allegro CL, available through Franz Inc. or
through third parties. Information on CORBA in general can be found at the Object Management Group's website, http://www.omg.org. By using CorbaLink, CL applications can
communicate in a seamless manner with applications written in other CORBA-supporting
languages such as C++, Java, Smalltalk, and others.

\section{TCP/IP Sockets}

TCP/IP provides a standard for communication over a network. Most Internet communication is done via TCP/IP. Both servers and clients can be written which use TCP/IP as the
underlying protocol.

In Allegro CL, the package \texttt{socket} provides CL-level functions which enable the use of
TCP/IP. As a simple example, we will show how to write both a client and a server which
can exchange data using TCP/IP sockets. These two programs can communicate when both
running on the same machine, or they can run on separate machines anywhere on a TCP/IP
network.

First we write the server\footnote{For space considerations we are not handling errors or trying to support multiple connections simultaneously.
See the Allegro CL documentation for socket examples which do provide full error-handling and simultaneous
connection support.}:

\begin{verbatim}
(in-package :user)

(defparameter *socket-server-port* 4444)

(defun open-server ()
  (socket:make-socket :connect :passive
                      :local-port *socket-server-port*))

(defun test-socket-server ()
  (let ((socket (open-server)))
    (unwind-protect
        (let ((stream (socket:accept-connection socket)))
          (let ((command (read stream)))
            (cond ((eql command :hello)
                   (format stream "Hello to you!~%")
                   (finish-output stream))
                  ((eql command :add-1-2)
                   (format stream "~a~%" (+ 1 2))
                   (finish-output stream))
                  (t (format stream "Sorry, I don't understand~%")
                     (finish-output stream))))
          (close stream))
      (close socket))))
\end{verbatim}

Now we write a simple client, which will contact the server, send a command, and print
out the response:

\begin{verbatim}
(in-package :user)

(defparameter *socket-server-port* 4444)

(defun test-socket-client (command)
  (let ((socket (socket:make-socket :remote-host "localhost"
                                    :remote-port *socket-server-port*)))
    (unwind-protect
        (progn
          (format socket "~s~%" command)
          (finish-output socket)
          (format t "~a~%" (read-line socket)))
      (close socket))))
\end{verbatim}

To test the server and client, start two CL processes. Load the server code into one process,
and call

\begin{verbatim}
(test-socket-server)
\end{verbatim}

in that process. Load the client code into the other process, and call

\begin{verbatim}
(test-socket-client :hello)
\end{verbatim}

in that process. In the client, you should see the response from the server printed out. In
the server, you should see the server return from the function call (since the server, as it is
written, only serves one request, then exits). Note that you could change ``localhost'' in the
definition of \texttt{test-socket-client} to any machine name on your network, and the client
and server could communicate with each other between different machines.

Also note that both the \texttt{socket:make-socket} and \texttt{socket:accept-connection} functions accept several optional keyword parameters, which are documented in Allegro CL's
online and hardcopy manuals.

\section{Databases}

Allegro CL provides connectivity to many SQL-based database systems, such as Oracle,
MySQL, PostgreSQL, Sybase and ODBC-compliant databases. Information on these packages, as well as more general information on accessing relational databases from CL, can be
found at

\begin{center}
\texttt{http://opensource.franz.com}
\end{center}

As a short example of database access from CL, we will present code for creating a table,
adding some data to it, and querying the data with MySQL (which is a free open-source
database server available on several platforms including Unix, Linux, and Windows).

\begin{verbatim}
(in-package :user)

(require :mysql)

(use-package :dbi.mysql)

(defun test-mysql ()
  (let ((db (connect "localhost" "testdb" "testuser" "testpasswd")))
    (sql db "create table captable (id int, name char(30))")
    (sql db "insert into captable values (~d, '~a')" 1 "Chicago")
    (sql db "insert into captable values (~d, '~a')" 2 "Austin")
    (sql db "insert into captable values (~d, '~a')" 3 "Philadelphia")
    (dolist (row (sql db "select * from captable where id > 1" :types :auto))
      (format t "~a: ~a~%" (first row) (second row)))
    (sql db "drop table captable")
    (disconnect db)))
\end{verbatim}

In the above example, ``testdb'' is the name of a database which must exist on the MySQL
server running on ``localhost,'' and ``testuser'' is a MySQL username which must have the
proper access privileges to the database (``testpasswd'' is the password).

Upon calling \texttt{(test-mysql)}, we expect to see the following output:

\begin{verbatim}
2: Austin
3: Philadelphia
\end{verbatim}

\section{Webserving}

Most CL vendors either provide direct support or provide third-party support for serving web
content through their CL environments. For Allegro CL, Franz Inc. provides AllegroServe
(``aserve''), a full-featured web application server. Applications built with aserve can provide
all the benefits of traditional web applications, with the additional advantage that all state
can be maintained in the running CL session. For traditional web applications, a new process
must typically be started for each connection; this leads to an inherent inefficiency since all
the application state must be fetched from a database or from persistent files every time.

AllegroServe can be obtained from

\begin{center}
\texttt{http://opensource.franz.com}
\end{center}

and is distributed under an open-source (LLGPL\footnote{LLGPL is the Lisp Lesser General Public License}) license.

\subsection{Publishing Static Pages}

To start the webserver, you use the function \texttt{net.aserve:start}. It takes several keyword
arguments. The only argument required is \texttt{:port}, which specifies the port number that
the server should listen on. The default HTTP port is 80. If you are running as a normal
(non-root) user on Unix/Linux, you will probably not have permission to run on port 80, so
you should choose a number above 1024 and note it for use in forming your URLs. ``8000''
is a traditional ``high port'' default for running webservers. The URL for your webserver
will then appear as \texttt{http://localhost:8000/} (of course replacing 8000 with whatever
port number you actually use).

In order to serve a simple plain HTML (i.e. static) page from a file, the function is
\texttt{publish-file}:

\begin{verbatim}
(net.aserve:start :port 8000)

(publish-file :path "/welcome.html" :file "/tmp/welcome.html")
\end{verbatim}

The \texttt{:path} argument gives the URL path, while the \texttt{:file} argument gives the name in
the Unix (or Windows) filesystem where the file actually resides.

To test the above example, create a simple HTML file and place it in \texttt{/tmp/welcome.html}
(or use the \texttt{:file} argument with a different location as appropriate). A simple HTML file
might look like this:

\begin{verbatim}
<html>
  <head>
    <title>Welcome</title>
  </head>
  <body>
    <h2>Welcome to my Web Page!</h2>
  </body>
</html>
\end{verbatim}

Now, start your web browser and visit the page \texttt{http://localhost:8000/welcome.html}.
You should see the web page come up in the browser.

\subsection{Publishing Dynamic Pages}

Dynamic web page serving is more interesting, since a function is called to generate the
HTML content on-the-fly. This function can tap into the entire running CL environment,
to query data structures, call external databases, etc. Here is an example:

\begin{verbatim}
(publish :path "/hello"
         :content-type "text/html"
         :function
         #'(lambda(req ent)
             (with-http-response (req ent)
               (with-http-body (req ent)
                 (html (:html
                        (:head (:title "Hello Page"))
                        (:body
                         ((:font :color "red")
                          (:h2 "Hello There!")))))))))
\end{verbatim}

To test the above example, visit the URL \texttt{http://localhost:8000/hello} in your browser.
You should see the heading ``Hello There!'' in red in your browser window.

You should also see in your CL console that some information is being logged on each
visit to the page.

Note the use of the \texttt{html} macro, which works somewhat like the \texttt{format}\index{format@\texttt{format}} function. It
takes a series of HTML ``tags'' represented as lists --- the first element of each list is a
keyword symbol representing the ``name'' of the tag, and the rest of the list are either
strings representing actual content, or more nested lists, or both.

\subsection{Query Parameters}

Query parameters are named values which can be spliced into a URL to pass data into
the page. Within AllegroServe they can be accessed with the \texttt{request-query-value}
function. As an example, consider the following \texttt{publish} form:

\begin{verbatim}
(publish :path "/welcome"
         :content-type "text/html"
         :function
         #'(lambda(req ent)
             (with-http-response (req ent)
               (with-http-body (req ent)
                 (let ((name (or (request-query-value "name" req)
                                 "Pal")))
                   (html (:html
                          (:head (:title "Welcome Page"))
                          (:body
                           ((:font :color "red")
                            (:h2 "Welcome, " (:princ name) "!"))))))))))
\end{verbatim}

To test this, visit \texttt{http://localhost:8000/welcome?name=Joe} in your browser. You should
see the heading ``Welcome, Joe!'' in the browser window.

\subsection{Form Parameters}

Query parameters are typically used for small amounts of data. For larger amounts of data,
you will probably want to use HTML form parameters. These are generally filled in by a
user through an interactive HTML form. To demonstrate the server side of such interaction,
we will use the \texttt{curl} Unix command to simulate the posting of form data to a published
aserve page.

Here is a \texttt{publish} form which can accept two form parameters named ``color'' and
``shape'':

\begin{verbatim}
(publish :path "/form-test"
         :content-type "text/html"
         :function
         #'(lambda(req ent)
             (with-http-response (req ent)
               (with-http-body (req ent)
                 (let ((color (request-query-value "color" req))
                       (shape (request-query-value "shape" req)))
                   (html (:html
                          (:head (:title "Form Test"))
                          (:body
                           (:h2 "You posted:")
                           (:princ (format nil "Color: ~a~%" color))
                           (:br)
                           (:princ (format nil "Shape: ~a" shape))))))))))
\end{verbatim}

Now, from a Unix terminal window, type the following command:

\begin{verbatim}
curl --data "color=red&shape=circle" http://localhost:8000/form-test
\end{verbatim}

You should see the following HTML printed on your console:

\begin{verbatim}
<html><head><title>Form Test</title></head>
<body><h2>You posted:</h2>Color: red<br />Shape: circle</body></html>
\end{verbatim}

Note that \texttt{request-query-value} can be used to retrieve both URL-encoded query parameters and form parameters --- aserve makes these available through the same interface.

\subsection{Multipart Form Data and File Uploads}

In order to handle binary files which are uploaded to the server, aserve provides the function
\texttt{get-multipart-content}. Combined with the AllegroServe documentation and/or the
aserve source code, this should be enough to get you started with handling uploaded files.

\subsection{Development and Production Modes}

AllegroServe has a concept of ``development'' and ``production'' mode which can be switched
with the \texttt{:debug} keyword argument. In debug mode, when an error is detected the handler
will stop with a break (see Section 2.8.3). This is usually convenient for development. In
production mode, on the other hand, aserve will simply output a standard HTTP error
response back to the browser and will not stop the server process.

The default debug setting for aserve is production mode. For development, you will
typically want to start aserve with:

\begin{verbatim}
(net.aserve:start :port 8000 :debug t)
\end{verbatim}

\subsection{Other Features}

This chapter has barely scratched the surface of what is available in AllegroServe. For a
complete treatment, see the documentation available at the AllegroServe website listed above.
Additional capabilities of note include:

\begin{itemize}
\item Authorization and security
\item Support for SSL (Secure Sockets Layer)
\item HTTP proxying
\item Virtual hosting (multiple hostnames and DNS names being served by the same physical
server)
\item Support for client-side HTTP from CL
\item Support for internationalization (international character sets) in URLs
\item Support for serving WebDAV (Web-based Distributed Authoring and Versioning)
\end{itemize}

\section{Embedded Languages}

This section applies to layered application development. In particular, we cover the concept
of an Embedded Language: a collection of macros and functions which add primitives and
other key operators to CL, effectively creating a ``language on top of CL.''

An example of this is GDL (Generative Design Language), which is the open-source
core of the Genworks\textregistered{} GDL commercial product. GDL is a general-purpose language
for Knowledge-Based Engineering, and is a significant superset of ANSI CL. Objects in
GDL are similar to CLOS\index{CLOS} objects, but add features such as cached demand-driven behavior,
automatic dependency tracking, and optional default values for unspecified data. GDL adds
new primitives to CL, such as the \texttt{define-object} operator, which is analogous to CLOS's
\texttt{defclass}\index{defclass@\texttt{defclass}}.

For layered applications and embedded languages, CL provides an environment unmatched
by other current languages. Because CL allows macros to generate other macros, and the
Lisp code can compute and generate Lisp code at all levels, it is straightforward to build
extremely high-level languages on top of CL. And because of the CL-based syntax, these
layered languages integrate seamlessly into the underlying CL environment.

\section{Conclusion}

Chapter 4 has provided a brief introduction into some of the interfacing capabilities available
in CL. For additional information on any of these topics, you should refer to the on-line
and printed documentation from Franz Inc., as well as to various resources available on the
Web.


% Chapter 5
\chapter{The Squeakymail Example}

This chapter will present a simple yet complete Common Lisp\index{Common Lisp} application. The application
is a web-based email client, called ``Squeakymail,'' which can be used for reading, composing,
sending, and managing email through a web browser. The complete code for this example
can be downloaded from

\begin{center}
\texttt{http://www.genworks.com/downloads}
\end{center}

In this chapter, we will give an overview of the application with excerpts from the source
code, to help tie together the concepts covered in earlier chapters. This example is intended
to give a flavor of working with CL on a real application.

\section{Requirements and Design Overview}

The basic requirements for this application are the following:

\begin{itemize}
\item Must be accessible from any web browser
\item Must handle POP3 mailboxes
\item Should allow the viewing of mail messages
\item Should allow the composition and sending of mail
\item Should allow the deletion and management of messages
\end{itemize}

Fortunately, Franz Inc. provides a free open-source library for handling POP3 email and
for composing and sending mail. This library can also be downloaded from the above URL.

The webserver will use AllegroServe, as introduced in Chapter 4. The overall design
breaks down as follows:

\begin{enumerate}
\item The AllegroServe session will store a list of messages in a global parameter
\item A background thread will periodically poll the POP3 server and update the list with
any new mail
\item The web interface will provide pages to view the list of messages, to read the content
of individual messages, and to compose and send new mail
\item Sufficient security will be in place to prevent multiple users of the application from
reading each other's email
\end{enumerate}

\section{Main Application File}

The main file for the application looks like this:

\begin{verbatim}
(in-package :user)

(defparameter *squeaky-port* 9876)

(defparameter *imap-server* "localhost")
(defparameter *pop-server* "localhost")
(defparameter *smtp-server* "localhost")

(defparameter *username* (sys:user-name))

(defparameter *poll-interval* 30)

(defun start-squeaky (&key (port *squeaky-port*))
  (start-server port)
  (start-background-polling))

(defun stop-squeaky ()
  (stop-background-polling)
  (stop-server))
\end{verbatim}

The function \texttt{start-squeaky} starts the webserver and kicks off a background thread to
poll the mailbox server for new mail. The \texttt{stop-squeaky} function stops this background
thread and shuts down the webserver.

\section{Working With Global State}

Because the web application runs in a persistent CL session, the entire program state can
be maintained in memory for fast and convenient access. No need to maintain state in a
database or in some transient or persistent file system.

For this application, we maintain the state in three main global variables:

\begin{verbatim}
(defparameter *lock* (mp:make-process-lock))

(defparameter *messages* nil)

(defparameter *pollingp* nil)
\end{verbatim}

The \texttt{*lock*} variable is a process lock, used for controlling access to the \texttt{*messages*} list
in a multi-threaded environment. The \texttt{*messages*} variable contains a list of mail-message
objects. The \texttt{*pollingp*} variable serves as a Boolean flag indicating whether the background
polling process is currently running.

\section{Background Polling}

The \texttt{start-background-polling} and \texttt{stop-background-polling} functions spawn and
halt a background thread which periodically polls for new mail:

\begin{verbatim}
(defun start-background-polling ()
  (setq *pollingp* t)
  (mp:process-run-function "Polling Process"
    #'polling-function))

(defun stop-background-polling ()
  (setq *pollingp* nil))

(defun polling-function ()
  (do () ((not *pollingp*))
    (update-messages-from-server)
    (sleep *poll-interval*)))
\end{verbatim}

The call to \texttt{mp:process-run-function} spawns a new operating system thread. This function takes a name for the thread, and a function which the thread should invoke. The
polling thread simply calls the \texttt{update-messages-from-server} function in a loop, waiting \texttt{*poll-interval*} seconds between each iteration.

The \texttt{update-messages-from-server} function is the function which actually contacts
the POP3 server:

\begin{verbatim}
(defun update-messages-from-server ()
  (mp:with-process-lock (*lock*)
    (let ((count (ignore-errors (pop:pop-message-count *pop-server*
                                                       *username*))))
      (when (and count (numberp count))
        (dotimes (n count)
          (unless (find n *messages* :key #'message-number)
            (let ((message (retrieve-message (1+ n))))
              (push message *messages*))))))))
\end{verbatim}

The \texttt{mp:with-process-lock} macro ensures that only one thread at a time can be accessing
and modifying the \texttt{*messages*} list.

The \texttt{pop:pop-message-count} function returns the number of messages currently waiting on the server for the given username. The username in this example defaults to the
Unix login name for the user running the CL session. You will have to ensure that you can
authenticate yourself to your POP3 server with this username.

If there are messages on the server which are not yet in the \texttt{*messages*} list, they are
retrieved and added to the list. The details of the \texttt{retrieve-message} function, which
wraps the calls to retrieve the actual message, are not shown.

\section{Serving the Web Pages}

The application provides three main web pages:

\begin{enumerate}
\item Index --- shows a list of the available messages, with links to read them
\item Read --- shows the content of an individual message
\item Compose --- allows the user to compose and send a new message
\end{enumerate}

Here is the function which publishes the main index page:

\begin{verbatim}
(publish :path "/squeaky"
         :content-type "text/html"
         :function
         #'(lambda (req ent)
             (with-http-response (req ent)
               (with-http-body (req ent)
                 (html
                  (:html
                   (:head (:title "Squeakymail"))
                   (:body
                    (:h1 "Your Mail")
                    ((:form :method "post" :action "/squeaky/compose")
                     ((:input :type "submit" :value "Compose New")))
                    (:hr)
                    (if *messages*
                        (html
                         ((:table :border 1)
                          (:tr (:th "Number") (:th "From")
                               (:th "Subject") (:th "Date"))
                          (dolist (message (reverse *messages*))
                            (html
                             (:tr
                              (:td ((:a :href
                                     (format nil "/squeaky/read?number=~a"
                                             (message-number message)))
                                    (:princ (message-number message))))
                              (:td (:princ (message-from message)))
                              (:td (:princ (message-subject message)))
                              (:td (:princ (message-date message))))))))
                        (html (:i "No messages")))
                    (:hr))))))))
\end{verbatim}

This function uses the \texttt{html} macro from AllegroServe to generate the HTML content
dynamically. The page includes a button to compose new mail, and a table showing all the
messages currently in the \texttt{*messages*} list. Each message number is a hyperlink to read that
particular message.

The ``read'' page is published in a similar manner:

\begin{verbatim}
(publish :path "/squeaky/read"
         :content-type "text/html"
         :function
         #'(lambda (req ent)
             (let* ((number-string (request-query-value "number" req))
                    (number (when number-string
                              (ignore-errors (parse-integer number-string))))
                    (message (when number
                               (find number *messages*
                                     :key #'message-number))))
               (with-http-response (req ent)
                 (with-http-body (req ent)
                   (if message
                       (html
                        (:html
                         (:head (:title "Read Mail"))
                         (:body
                          (:h2 "Message " (:princ number))
                          (:b "From: ") (:princ (message-from message)) (:br)
                          (:b "Subject: ") (:princ (message-subject message))
                          (:br)
                          (:b "Date: ") (:princ (message-date message)) (:br)
                          (:hr)
                          (:pre (:princ (message-body message)))
                          (:hr)
                          ((:form :method "post" :action "/squeaky/delete")
                           ((:input :type "hidden" :name "number"
                                    :value (format nil "~a" number)))
                           ((:input :type "submit" :value "Delete")))
                          ((:a :href "/squeaky") "Back to Index"))))
                       (html
                        (:html
                         (:head (:title "Error"))
                         (:body
                          (:h2 "Message not found")
                          ((:a :href "/squeaky") "Back to Index"))))))))))
\end{verbatim}

This page retrieves the message number from the query string, finds the corresponding
message in the \texttt{*messages*} list, and displays its content. It also provides a button to
delete the message and a link back to the main index.

\section{Composing and Sending Mail}

The compose page is slightly more complex, as it needs to handle both displaying a blank
form and processing the submitted form data:

\begin{verbatim}
(publish :path "/squeaky/compose"
         :content-type "text/html"
         :function
         #'(lambda (req ent)
             (let ((to (request-query-value "to" req))
                   (subject (request-query-value "subject" req))
                   (body (request-query-value "body" req)))
               (with-http-response (req ent)
                 (with-http-body (req ent)
                   (if (and to subject body)
                       (progn
                         (send-smtp-mail *smtp-server* *username*
                                        to subject body)
                         (html
                          (:html
                           (:head (:title "Message Sent"))
                           (:body
                            (:h2 "Your message has been sent")
                            ((:a :href "/squeaky") "Back to Index")))))
                       (html
                        (:html
                         (:head (:title "Compose Mail"))
                         (:body
                          (:h2 "Compose New Message")
                          ((:form :method "post" :action "/squeaky/compose")
                           (:b "To: ")
                           ((:input :type "text" :name "to" :size 40)) (:br)
                           (:b "Subject: ")
                           ((:input :type "text" :name "subject" :size 40))
                           (:br) (:br)
                           ((:textarea :name "body" :rows 20 :cols 60))
                           (:br) (:br)
                           ((:input :type "submit" :value "Send"))
                           "&nbsp;&nbsp;"
                           ((:input :type "button" :value "Cancel"
                                    :onclick "window.location='/squeaky'")))
                          (:hr)))))))))))
\end{verbatim}

If the form has been submitted (i.e., if the ``to,'' ``subject,'' and ``body'' parameters are present),
the \texttt{send-smtp-mail} function is called to send the message. Otherwise, a blank form is
displayed.

\section{Deleting Messages}

The delete function removes a message from the local list and also deletes it from the POP3
server:

\begin{verbatim}
(publish :path "/squeaky/delete"
         :content-type "text/html"
         :function
         #'(lambda (req ent)
             (let* ((number-string (request-query-value "number" req))
                    (number (when number-string
                              (ignore-errors (parse-integer number-string)))))
               (when number
                 (mp:with-process-lock (*lock*)
                   (setq *messages* (remove number *messages*
                                           :key #'message-number))
                   (ignore-errors
                     (pop:delete-letter *pop-server* *username* number))))
               (with-http-response (req ent)
                 (with-http-body (req ent)
                   (html
                    (:html
                     (:head (:title "Message Deleted"))
                     (:body
                      (:h2 "Message has been deleted")
                      ((:a :href "/squeaky") "Back to Index")))))))))
\end{verbatim}

Note the use of \texttt{mp:with-process-lock} to ensure thread-safe modification of the \texttt{*messages*}
list.

\section{Starting the Webserver}

The \texttt{start-server} and \texttt{stop-server} functions are simple wrappers around the AllegroServe functions:

\begin{verbatim}
(defun start-server (port)
  (net.aserve:start :port port))

(defun stop-server ()
  (net.aserve:shutdown))
\end{verbatim}

\section{Message Data Structure}

The application uses a simple structure to represent mail messages:

\begin{verbatim}
(defstruct message
  number
  from
  subject
  date
  body)
\end{verbatim}

The \texttt{defstruct}\index{defstruct@\texttt{defstruct}} macro automatically creates constructor, accessor, and type-predicate functions for this structure.

\section{Testing the Application}

To test Squeakymail, you will need:

\begin{enumerate}
\item A running POP3 mail server (or access to one)
\item A running SMTP mail server (or access to one)
\item The Post Office library from Franz Inc.
\item AllegroServe
\end{enumerate}

Once these components are in place, you can start the application by loading all the source
files and calling:

\begin{verbatim}
(start-squeaky)
\end{verbatim}

Then point your web browser to:

\begin{verbatim}
http://localhost:9876/squeaky
\end{verbatim}

You should see the main index page with your mail messages listed. From there you can
read, compose, and delete messages.

\section{Security Considerations}

This example application is intended for demonstration purposes only and should not be used
in a production environment without significant security enhancements. In particular:

\begin{itemize}
\item The application does not implement user authentication
\item Passwords are not encrypted
\item There is no session management to prevent unauthorized access
\item The application assumes a single user
\end{itemize}

For a production application, you would need to implement proper authentication, session
management, and access controls.

\section{Enhancements}

Some possible enhancements to this application might include:

\begin{itemize}
\item Support for IMAP in addition to POP3
\item Ability to reply to messages
\item Support for attachments
\item Better formatting of HTML email
\item Address book functionality
\item Folders and message filtering
\item Search functionality
\item Support for multiple users with authentication
\end{itemize}

\section{Conclusion}

This chapter has demonstrated a complete, albeit simple, web application written in Common
Lisp. The example shows how the various concepts covered in earlier chapters come together
in a real application:

\begin{itemize}
\item Use of global parameters to maintain state
\item Multi-threaded programming with process locks
\item Web serving with AllegroServe
\item Dynamic HTML generation
\item External library integration (POP3/SMTP)
\item Data structures and list manipulation
\end{itemize}

The complete source code for Squeakymail, along with installation instructions, can be
downloaded from the URL given at the beginning of this chapter. We encourage you to
download, install, and experiment with this application as a learning tool for Common Lisp
development.


\backmatter

\printindex

\end{document}
